\section{Correspondences, tautological classes, and extension problems}
\label{sec:part-c}
\providecommand{\Lie}{\operatorname{Lie}}
\providecommand{\tr}{\operatorname{tr}}
\providecommand{\Alt}{\operatorname{Alt}}
\providecommand{\sgn}{\operatorname{sgn}}
\providecommand{\Gr}{\operatorname{Gr}}
\providecommand{\id}{\operatorname{id}}

\subsection{C.1: a finite presentation for every cohomological tautological ring}

We use a Weil cohomology with coefficient field $L$ of characteristic zero, the usual Koszul commutativity constraint, and cohomological grading. Let $X$ be smooth, projective, geometrically connected, and of dimension $n>0$ over an algebraically closed field. Put $D=2n$, $A=H^*(X)$, and $\kappa(a)=\int_Xa$. We choose the top class $y$ with $\kappa(y)=1$. All the generators under consideration have degree divisible by $D$; we write $T_k[d]=T_k\cap H^{Dd}(X^k)$. For $n=0$, $X$ is a point and $T_k=L$ in degree zero, which we henceforth set aside.

\begin{definition}[Pointed partitions]
\label{c:def-pointed}
\status{Definition, with degrees verified}
Let $\mathcal P_k$ be the set of partitions of $\{0,1,\ldots,k\}$. The block containing $0$ is the distinguished block. For $P\in\mathcal P_k$, put $J=P(0)\setminus\{0\}$ and define
\[
 \varepsilon_P=\prod_{i\in J}y_i\prod_{B\in P,\,0\notin B}\delta_B,
 \qquad d(P)=k-|P|+1.
\]
Here $\delta_B$ is the pullback of the small diagonal in $X^B$ and $\delta_{\{i\}}=1$. Thus $\varepsilon_P\in T_k[d(P)]$. This describes the pointed and unpointed partial diagonals of manuscript 11:55 without duplicate notation.
\end{definition}

\begin{theorem}[C.1(b): the complete graph product rule]
\label{c:graph-rule}
\status{Proved}
Choose a tree on each block of $P$ and of $Q$, using the edges $\{0,i\}$ in a distinguished block, and take the union as a multigraph on $\{0,\ldots,k\}$. Write $b(C)=|E(C)|-|V(C)|+1$ for the cycle rank of a connected component. A component not containing $0$ contributes
\[
 \begin{cases}
 \delta_C,&b(C)=0,\\
 \chi(X)\prod_{i\in C}y_i,&b(C)=1,\\
 0,&b(C)\geq2,
 \end{cases}
\]
whereas the component containing $0$ contributes $\prod_{i\in C\setminus\{0\}}y_i$ if $b(C)=0$, and zero otherwise. Multiplying these contributions gives $\varepsilon_P\varepsilon_Q$, independently of the chosen trees. In particular, the $\varepsilon_P$ span $T_k$.
\end{theorem}

\begin{proof}
A tree of pairwise equalities imposes a small diagonal, since its edges impose independent diagonal conditions. In the algebra $A^{\otimes k}$ this also follows by composing the coevaluation for the perfect pairing $\kappa(ab)$. After the tree conditions have been imposed, each further edge in the same component is the self-intersection of the diagonal. Its pullback is the Euler element
\[
 e_A=\mu(\Delta_A)=\chi(X)y,
\]
where $\Delta_A\in A\otimes A$ is the coevaluation and $\mu$ is multiplication. To verify the last equality directly, choose homogeneous dual bases; the coefficient of $y$ is the supertrace of the identity, namely $\sum_r(-1)^r\dim H^r(X)$. Since $n>0$, $y^2=0$. Consequently a second excess edge gives zero. If the component already contains a point condition, an excess edge gives $y e_A=0$; the tree itself makes every vertex in that component equal to the point. The coevaluation identities show that changing the tree leaves the same diagonal class. Every product of generators is therefore either zero or a scalar times a pointed-partition class. These arguments establish the previously unchecked products of 11:55--57, including the cycle and point cases.
\end{proof}

We next give an actual finite system of generators for all relations, valid even when the restricted intersection pairing on $T_k$ is degenerate. Using only that restricted pairing would not suffice.

\begin{theorem}[C.1(a): all relations from the multilinear forms]
\label{c:matrix-presentation}
\status{Proved}
Choose a homogeneous basis $(e_a)$ of $A$, containing $e_0=1$. For fixed $d$, form a matrix $C_{k,d}$ whose columns are the partitions $P$ with $d(P)=d$ and whose rows are the basis tensors $e_{a_1}\otimes\cdots\otimes e_{a_k}$ of total degree $D(k-d)$. Its entry is
\begin{equation}
\label{c:contraction-matrix}
 (C_{k,d})_{\mathbf a,P}
 =\epsilon(\mathbf a,P)
 \prod_{B\in P,\,0\notin B}\kappa\Bigl(\prod_{i\in B}e_{a_i}\Bigr)
 \prod_{i\in P(0)\setminus\{0\}}[a_i=0].
\end{equation}
In each block the indices are in increasing order; $\epsilon$ is the Koszul sign of collecting the factors into those blocks. Then
\[
 \dim_L T_k[d]=\operatorname{rank}_L C_{k,d},\qquad
 \sum_Pc_P\varepsilon_P=0\ \Longleftrightarrow\ (c_P)_P\in\ker C_{k,d}.
\]
A finite presentation of $T_k$ consists of one variable $z_P$ for each $P\in\mathcal P_k$, the unit relation, all product relations of Theorem~\ref{c:graph-rule}, and one linear relation for each vector in a basis of each $\ker C_{k,d}$. Thus Gaussian elimination gives generators for the entire ideal of relations. The input is exactly the graded vector space, its unit and trace, and the multilinear forms $\kappa(u_1\cdots u_m)$ on positive-degree elements.
\end{theorem}

\begin{proof}
The projection formula for each diagonal gives
\[
 \int_{X^k}\varepsilon_P(e_{a_1}\otimes\cdots\otimes e_{a_k})
 =(C_{k,d})_{\mathbf a,P}.
\]
For a pinned coordinate, $\kappa(ye_a)=[a=0]$; for an unpinned block it is the integral of the product of its inputs. These are all factors, with precisely the displayed Koszul permutation sign. The full Poincar\'e pairing on $A^{\otimes k}$ is perfect, so its contraction with all rows detects the zero vector. This proves both the kernel and rank assertions. The graph relations reduce any polynomial to a linear combination of $z_P$; the matrix kernel is exactly the kernel remaining after this reduction. No nondegeneracy of the restricted pairing on the span of the $\varepsilon_P$ is used. Finally, multilinear forms involving a unit are obtained by deleting that unit, so the stated positive-degree data suffice.
\end{proof}

\begin{proposition}[The modified diagonal and the correction to R8]
\label{c:modified-diagonal}
\status{Proved; proposed generation by this relation alone refuted}
The class
\[
 M=\delta_{123}-\delta_{12}y_3-\delta_{23}y_1-\delta_{13}y_2
       +y_1y_2+y_2y_3+y_1y_3
 \in H^{4n}(X^3)
\]
vanishes if and only if $\kappa(uvw)=0$ for every three positive-degree classes $u,v,w$. The degree $H^{2n}$ stated in R8 is incorrect when $n$ denotes the geometric dimension. Neither $M$ nor its higher-arity analogues, without further dimension relations, generate all relations in every $T_k$.
\end{proposition}

\begin{proof}
Pair $M$ with $u\otimes v\otimes w$. If all three inputs have positive degree, all terms except the small diagonal vanish and the pairing is $\kappa(uvw)$. If at least one input is $1$, the terms cancel by inclusion and exclusion; multilinearity reduces every input to these cases. Poincar\'e duality proves the equivalence. For the generation assertion take $X=\mathbb P^1$ and $k=2$. There is the linear relation $\delta_{12}=y_1+y_2$; relations with at least three distinct indices are absent in this ring, while the universal self-intersection and point relations have polynomial degree at least two. Their homogeneous ideal cannot contain this nonzero linear polynomial. More generally, the dimension relations in Theorem~\ref{c:pair-ranks} below are additional indispensable relations. We distinguish this correction from the valid triple-product criterion in R8.
\end{proof}

\begin{theorem}[Dimensions for curves and for surfaces with $b_1=0$]
\label{c:pair-ranks}
\status{Proved}
For a curve of genus $g$, put $b=-2g$; for a smooth projective surface with $b_1=0$, put $b=b_2$. Define $q_{ij}=\delta_{ij}-y_i-y_j$. All relations are generated by
\begin{align}
 y_i^2&=0,&y_iq_{ij}&=0,&q_{ij}^2&=b\,y_iy_j,&q_{ij}q_{ik}&=y_iq_{jk},\label{c:pair-local}
\end{align}
and the following dimension relations, with all displayed indices distinct. For a surface they are
\begin{equation}
\label{c:orth-det}
 \sum_{\sigma\in S_{b+1}}\sgn(\sigma)
   \prod_{a=1}^{b+1}q_{i_a,j_{\sigma(a)}}=0;
\end{equation}
for a curve they are
\begin{equation}
\label{c:sympl-pf}
 \sum_{\text{perfect matchings }\mathcal M\text{ on }2g+2\text{ slots}}
       \prod_{\{i,j\}\in\mathcal M}q_{ij}=0.
\end{equation}
There is no sign in \eqref{c:sympl-pf}, because the symplectic signs are already included by the odd degree of $H^1$ and the Koszul constraint.

For a partition $\lambda$, write $f^\lambda$ for the number of standard Young tableaux of shape $\lambda$, explicitly $|\lambda|!$ divided by the product of the hook lengths. Put
\[
 r_b(m)=
 \begin{cases}
 \displaystyle\sum_{\lambda\vdash m,\,\ell(\lambda)\leq b} f^{2\lambda},&b\geq0\quad\text{(surface)},\\[4pt]
 \displaystyle\sum_{\lambda\vdash m,\,\lambda_1\leq g} f^{2\lambda},&b=-2g\quad\text{(curve)},
 \end{cases}
 \qquad r_b(0)=1.
\]
Then the answer in every degree is
\begin{equation}
\label{c:dimension-pair}
 \dim T_k[d]=
 \sum_{\substack{m\geq0,\ a=d-m\geq0\\a+2m\leq k}}
 \frac{k!}{a!(2m)!(k-a-2m)!}\,r_b(m).
\end{equation}
The curve formula includes $g=0$, when $r_0(m)=0$ for $m>0$ and the answer is $\binom{k}{d}$.
\end{theorem}

\begin{proof}
In both cases $A=L1\oplus W\oplus Ly$, with $W=H^1$ for a curve and $W=H^2$ for a surface. The product on $W$ is a perfect alternating, respectively symmetric, pairing with values in $Ly$. Direct contraction gives \eqref{c:pair-local}; the superdimension of $W$ gives $b$. In particular the positive triple products vanish. Repeated use of those relations reduces a monomial to $y_I\prod_{\{i,j\}\in\mathcal M}q_{ij}$, with disjoint indices. Factorwise cohomological degrees separate different point sets $I$ and different sets of paired vertices. For a fixed set of $2m$ paired vertices, the remaining problem is exactly the span of the pairing tensors in $W^{\otimes2m}$.

For completeness, the matching representation of $S_{2m}$ decomposes without multiplicity into the Specht modules indexed by the even-row shapes $2\lambda$. Its contraction matrix has nonzero eigenvalue exactly for $\ell(\lambda)\leq b$ in the orthogonal case and $\lambda_1\leq g$ in the symplectic case. The orthogonal calculation and the argument that every kernel Specht module is generated by the first determinant are stated and proved in \citet[Proposition 3.5 and \S3.6, pp.~508--510]{c_yin2015}; the symplectic calculation, including the all-plus convention and generation by \eqref{c:sympl-pf}, is worked out in \citet[\S6, equation (2), and the argument preceding Corollary 6.4]{c_tavakol2018}. These calculations concern a matrix of pairings and therefore apply to the present $W$ independently of whether it comes from a K3 surface or a hyperelliptic curve. To see the generation step, a kernel Young diagram has either a column of length $b+1$ or a row of length $2g+2$; antisymmetrizing that column, respectively symmetrizing that row with the Koszul signs, supplies \eqref{c:orth-det}, respectively \eqref{c:sympl-pf}. Decomposing the larger stabilizer into cosets reduces every larger kernel generator to its pullbacks and products. The hook formula gives $r_b(m)$. Finally, choose the $a$ point slots and the $2m$ pairing slots; the number of choices is the multinomial coefficient in \eqref{c:dimension-pair}. Summing gives the formula and proves completeness of the relations.
\end{proof}

\begin{proposition}[Explicit dimension algorithms for projective spaces and abelian varieties]
\label{c:other-dimensions}
\status{Proved}
The following finite integer matrices determine every dimension by ordinary exact rational rank; in particular they give a uniform computable answer for every $n,g,k,d$, with no geometric unknown in their entries.

For $X=\mathbb P^n$, the columns are $P\in\mathcal P_k$ with $d(P)=d$, and the rows are $(a_1,\ldots,a_k)\in\{0,\ldots,n\}^k$ with $\sum a_i=nd$. The entry is $1$ precisely when
\[
 a_i=n\quad(i\in P(0)\setminus\{0\}),\qquad
 \sum_{i\in B}a_i=n(|B|-1)\quad(0\notin B),
\]
and is zero otherwise. Its rank is $\dim T_k[d]$.

For an abelian variety of dimension $g>0$, use the exterior algebra on the $2gk$ symbols $e_{i,a}$, ordered lexicographically. In this algebra put
\[
 y_i=\prod_{a=1}^{2g}e_{i,a},\qquad
 \delta_{ij}=\prod_{a=1}^{2g}(e_{i,a}-e_{j,a}).
\]
For each $P$ multiply these expressions along any of its trees and its pinned vertices. The coefficient matrix in the exterior-monomial basis of degree $2gd$ has entries in $\Z$ and rank $\dim T_k[d]$. These ranks depend only on $g,k,d$.
\end{proposition}

\begin{proof}
For projective space, $H^*(X)=L[h]/(h^{n+1})$, $\int h^n=1$, and the small diagonal on a block $B$ is the sum of the monomials satisfying the displayed condition, each with coefficient one. This follows either by pairing with monomials or by induction along a tree. For an abelian variety, $H^*(X)=\bigwedge H^1(X)$. Choose a volume-normalized basis; the diagonal formula follows from the difference map $(x,y)\mapsto x-y$, whose pullback on $H^1$ is $e\mapsto e\otimes1-1\otimes e$, and whose inverse image of the origin is the diagonal. Expansion in the basis is therefore an explicit expression for every $\varepsilon_P$. In both cases the coefficient matrix has precisely these classes as columns, so its rank is their dimension. These matrices are distinct from an unevaluated Chow-theoretic condition: their entries require only finite integer arithmetic.
\end{proof}

The exact-arithmetic program \texttt{code/part\_c/check\_tautological.py} implements the Frobenius algebra construction independently of the display formulas. It verified 2,286 graph products and the curve/surface rank formula in the finite ranges recorded in its JSON output. For example, the dimensions for $k=4$, in normalized degrees $d=0,\ldots,4$, are
\[
\begin{array}{c|ccccc}
 X&0&1&2&3&4\\\hline
 \mathbb P^1&1&4&6&4&1\\
 \mathbb P^2&1&10&19&10&1\\
 \mathbb P^3&1&10&25&14&1\\
 \text{curve }g=1&1&10&20&10&1\\
 \text{curve }g=2&1&10&21&10&1
\end{array}
\]
The non-palindromic row for $\mathbb P^3$ is one reason not to replace the full pairing in Theorem~\ref{c:matrix-presentation} by an assumed perfect tautological pairing.

\subsection{C.1(c): Chow lifting, and the exact K3 obstruction}

Write $T_k^{\CH}(X,o)$ for the rational Chow subring generated by the same diagonals and a degree-one zero-cycle $o$. The requested lifting of all cohomological relations is exactly the injectivity of the cycle class map on this subring. In Chow groups the graph proof has Euler class $c_n(T_X)$ in place of $\chi(X)o$; this replacement itself is a condition to verify, not a consequence of the cohomological computation.

\begin{theorem}[The abelian and hyperelliptic cases]
\label{c:chow-positive}
\status{Known, with the required implication proved}
For an abelian variety with $o$ its origin, every cohomological relation in the rings of C.1(a) holds in $T_k^{\CH}(X,o)$ for every $k$. The same conclusion holds for a hyperelliptic curve with $o$ a Weierstrass point (and for an elliptic curve with its origin and for $\mathbb P^1$).
\end{theorem}

\begin{proof}
For an abelian variety, \citet[Theorem 6.2.5 and Corollary 6.2.6]{c_osullivan2011} identifies the algebra of symmetrically distinguished cycles with the numerical Chow algebra, compatibly with products and with pullbacks and pushforwards by homomorphisms. The origin is the pushforward of $1$ by the zero homomorphism from the zero-dimensional abelian variety. Each diagonal is the pushforward of $1$ by a diagonal homomorphism. Their products and pullbacks are consequently symmetrically distinguished. A homologically zero cycle is numerically zero; the stated isomorphism then forces a distinguished such cycle to vanish in rational Chow groups. This establishes the assertion for precisely the requested generators, without assuming the Hodge or Tate conjecture.

For a hyperelliptic curve of genus $g\geq2$, $K_X=(2g-2)o$ in $\Pic(X)\otimes\Q$. The modified diagonal vanishes by \citet[Proposition 4.8 and its relation with the cycle on $X^3$ in \S4]{c_grossschoen1995}; equivalently one may use the explicit divisor calculation in that proposition followed by the finite symmetric-power push--pull. The remaining relations are precisely \eqref{c:pair-local} and \eqref{c:sympl-pf}. \citet[\S6, equation (2), and Corollary 6.4]{c_tavakol2018} proves them and the perfectness of every tautological intersection pairing on the fixed curve. A homologically zero tautological cycle pairs to zero with every complementary tautological cycle, hence is zero by that perfectness. The identification of Tavakol's class $a_i=K_i/(2g-2)$ with $o_i$ verifies the change of generators. Genus one is covered both by the abelian argument and by Tavakol's Example 6.2. Genus zero is the polynomial Chow presentation of projective space.
\end{proof}

\begin{theorem}[Classification of the allowed zero-cycles on an abelian variety]
\label{c:abelian-point-classification}
\status{Proved over every algebraically closed field}
Let $A/F$ be an abelian variety over an algebraically closed field of any characteristic, and let $o\in\CH_0(A)_\Q$ have degree one. The following are equivalent:
\begin{enumerate}[label=(\roman*)]
\item all the cohomological point-and-diagonal relations hold in $T_k^{\CH}(A,o)$ for every $k$;
\item $\Delta_*o=o\boxtimes o$ in $\CH_0(A\times A)_\Q$;
\item $o=[a]$ in $\CH_0(A)_\Q$ for a point $a\in A(F)$.
\end{enumerate}
The point is unique modulo torsion in $A(F)$. In particular, specifying the origin is sufficient but is not necessary for the requested lifting.
\end{theorem}

\begin{proof}
The equality in (ii) is the cohomological relation $\delta_{12}y_1=y_1y_2$, since both classes are degree-one zero-cycles on $A^2$. Thus (i) implies (ii). If (iii) holds, simultaneous translation by $a$ on every factor preserves all partial diagonals and carries the origin class to $[a]$. Theorem~\ref{c:chow-positive} proves (i).

We prove the remaining implication algebraically, so that no characteristic-zero restriction is hidden. Put $g=\dim A$ and write $\widehat A$ for the dual abelian variety. For the normalized Poincar\'e bundle $\mathcal P$ on $A\times\widehat A$, define the rational Chow Fourier transform by
\[
 \mathcal F_A(z)=(p_2)_*\bigl(p_1^*z\,\exp(c_1(\mathcal P))\bigr).
\]
It is an isomorphism and $\mathcal F_{\widehat A}\mathcal F_A=(-1)^g[-1]^*$. These identities hold over the arbitrary algebraically closed field of \citet[Theorem 2.2, p.~156, with the field convention on p.~154]{c_mukai1981}: apply Grothendieck--Riemann--Roch to his Fourier--Mukai identity, using the trivial tangent bundle of an abelian variety. The Chern character identifies rational $K$-theory with rational Chow groups. We will also use the directly checked Poincar\'e identities below; over $\C$ they are formulas F1--F3 and Proposition 1 of \citet[pp.~647--648]{c_beauville1986}.

Set $z=\mathcal F_A(o)$. Its codimension-zero part is $\deg(o)=1$. Fourier transformation of (ii), using the defining Poincar\'e bundle and the projection formula, gives
\[
 m_{\widehat A}^*z=p_1^*z\,p_2^*z.
\]
Indeed Fourier transformation changes $\Delta_*$ into pullback by the dual homomorphism $m_{\widehat A}$, and changes an exterior product into the exterior product of the transforms. The ideal of positive-codimension classes is nilpotent, so the finite series $D=\log z$ is defined and satisfies
\[
 m_{\widehat A}^*D=p_1^*D+p_2^*D,
 \qquad [n]^*D=nD\quad(n\geq1).
\]

Write $D=\sum_{r=1}^gD_r$ by codimension. We verify that $D_r=0$ for $r\geq2$ without assuming the conjectural vanishing of any negative Beauville piece. Choose an integer $n\geq2$ prime to the characteristic. If $w_q$ is the codimension-$q$ part of $\mathcal F_{\widehat A}(D_r)$, its defining formula uses the power $c_1(\mathcal P)^{g-r+q}$. The identity $(1\times[n])^*\mathcal P=\mathcal P^{\otimes n}$ and flat base change therefore give
\[
 [n]^*w_q=n^{g-r+q}w_q.
\]
Fourier functoriality for the isogeny $[n]$, applied to $[n]^*D_r=nD_r$, gives $[n]_*w_q=nw_q$. Since $[n]_*[n]^*=n^{2g}$, a nonzero $w_q$ requires
\[
 n^{2g}=n^{g-r+q+1},\qquad q=g+r-1.
\]
The inequality $q\leq g$ forces $r\leq1$. Fourier inversion now gives $D=D_1$. On the N\'eron--Severi group $[n]^*$ acts by $n^2$, whereas it acts on $D_1$ by $n$. Hence $D_1\in\Pic^0(\widehat A)\otimes\Q$.

Biduality identifies $\Pic^0(\widehat A)$ with $A(F)$. This group is divisible: multiplication by every positive integer is a surjective isogeny and is surjective on points over an algebraically closed field, also when the integer is divisible by the characteristic. Thus there is $a\in A(F)$ with $D_1=c_1(\mathcal P_a)$ in rational Chow groups. Directly from its definition, $\mathcal F_A([a])=\exp(c_1(\mathcal P_a))=z$. Fourier inversion proves $o=[a]$. Finally $[a]=[b]$ in rational Chow groups if and only if their Albanese classes differ by torsion: the forward implication follows from Albanese, and the reverse from the identical Fourier exponentials of torsion-equivalent points. This proves the uniqueness assertion and the theorem.
\end{proof}

\begin{proposition}[A finite criterion for curves]
\label{c:curve-chow-criterion}
\status{Proved over every algebraically closed field}
Let $C$ be a smooth projective connected curve and $o\in\CH_0(C)_\Q$ have degree one. All point-and-diagonal cohomological relations hold in all its powers if and only if $M_o=0$ in $\CH^2(C^3)_\Q$. If $g\geq2$, this forces
\[
 o=\frac{K_C}{2g-2}\quad\text{in }\CH^1(C)_\Q.
\]
For $g=0$ or $1$ every degree-one rational zero-cycle has the required property. Over $\C$, the genus-at-least-two criterion is also the multiplicativity criterion in \citet[Propositions 3.1 and 3.5]{c_fulaterveervial2021}.
\end{proposition}

\begin{proof}
Necessity follows because $M_o$ is always cohomologically zero for a curve. Suppose it vanishes. Intersecting with $p_1^*o$ and pushing to the last two factors gives
\[
 \Delta_*o=o\boxtimes o.
\]
Intersecting instead with $\delta_{12}$ and pushing to the third factor gives
\[
 c_1(T_C)=(2-2g)o.
\]
In the second calculation, the small-diagonal term contributes $c_1(T_C)$, the term $\delta_{12}o_3$ contributes $(2-2g)o$, and the four nonzero remaining point terms cancel in pairs. These identities and the expansion of $M_o=0$ give all the local relations \eqref{c:pair-local} for $q_{ij}=\delta_{ij}-o_i-o_j$, with parameter $b=-2g$.

We check the remaining matching relation, including for a rational degree-one zero-cycle. Choose a base point $e$ and the Abel--Jacobi map $C\to J(C)$. The Jacobian's principal autoduality and the theorem of the square identify the pullback of its normalized symmetric Poincar\'e class, up to a common sign, with $\delta_{12}-e_1-e_2$. Translation of the Abel--Jacobi map by $z\in J(C)$ adds the corresponding algebraically trivial divisor to each factor. Since $J(C)(k)$ is divisible, we may choose $z$ so that the new pullback is $q_{12}$ in rational Chow groups, up to the same sign. Thus every matching monomial in the $q_{ij}$ is the pullback of a matching monomial in Poincar\'e divisors on a power of $J(C)$.

Let $\theta$ be a symmetric principal polarization on $J(C)$ and use integer parameters $n_1,\ldots,n_{2g+2}$. Pull back $\theta^{g+1}=0$ along the homomorphism $(x_i)\mapsto\sum n_ix_i$. The theorem of the square gives a polynomial identity whose quadratic form is
\[
 \sum_i n_i^2\theta_i+\sum_{i<j}n_in_jP_{ij}.
\]
The coefficient of $n_1\cdots n_{2g+2}$ in its $(g+1)$st power is $(g+1)!$ times the sum over all perfect matchings of the products of $P_{ij}$. A rational polynomial vanishing at every integer tuple has every coefficient zero. Pulling back to $C^{2g+2}$ proves precisely \eqref{c:sympl-pf}. Theorem~\ref{c:pair-ranks} now proves sufficiency in every power. These are algebraic arguments valid also in characteristic two.

For genus zero, $\CH_0(C)_\Q=\Q$. For genus one, choose an origin. The degree-zero Chow group is its Jacobian group, so divisibility shows that every degree-one rational zero-cycle equals a point class. Apply Theorem~\ref{c:abelian-point-classification}. The formula for $o$ when $g\geq2$ follows from the second displayed identity.
\end{proof}

\begin{proposition}[The exact additional obstruction for regular surfaces]
\label{c:surface-chow-criterion}
\status{Proved equivalence; does not assume finite-dimensionality}
Let $S/\C$ be a smooth projective surface with $b_1(S)=0$, and let $o$ have degree one. All the point-and-diagonal cohomological relations lift in all powers if and only if $M_o=0$ and $h(S)$ is Kimura finite-dimensional. If they hold, $o$ is uniquely determined: for every ample divisor $H$,
\[
 o=\frac{H^2}{\deg H^2}.
\]
\end{proposition}

\begin{proof}
The projector $\pi_2=\Delta-o\times S-S\times o$ has realization $H^2(S)$; the other two summands are Tate. Write $b=b_2(S)$. The alternating cross-matching determinant on $2(b+1)$ factors is $(b+1)!$ times the alternating projector on $h^2(S)^{\otimes(b+1)}$. Therefore the lifting of all relations implies $\bigwedge^{b+1}h^2(S)=0$, hence finite-dimensionality. It also implies $M_o=0$, since three positive-degree classes have zero product on such a surface.

Conversely, let $M_o=0$. Acting on $o$, on the diagonal, and on $H\boxtimes H$ gives respectively
\[
 \Delta_*o=o\boxtimes o,\qquad c_2(T_S)=(b+2)o,
 \qquad H^2=(\deg H^2)o.
\]
The same expansions as for the curve give all local relations \eqref{c:pair-local}, now with $b=b_2(S)$. Finite-dimensionality supplies the remaining alternating determinant: the realization is purely even of dimension $b$, so its Kimura exterior dimension is $b$. Theorem~\ref{c:pair-ranks} says that the local identities and these determinants generate every cohomological relation. This proves sufficiency and uniqueness. This formal argument is consistent with the larger divisor-generated criterion in \citet[Proposition 4.5]{c_fulaterveervial2021}; it explicitly applies to the narrower ring asked for here.
\end{proof}

\begin{theorem}[A finite cycle criterion when the cohomology is even]
\label{c:finite-even-criterion}
\status{Proved in the pure-even case}
Let $X/\C$ be smooth projective and connected, with $H^{\rm odd}(X,\Q)=0$,
let $n=\dim X>0$, and let $o\in\CH^n(X)_\Q$ have degree $1$.
Put $d=\dim_\Q H^*(X,\Q)=\chi(X)$.  There is a finite list of rational
tensor expressions $P_1,\ldots,P_s$, constructed by the procedure below
from the rational cup-product table and trace of $X$, for which the
following criterion is necessary and sufficient for all the relations
of Theorem~\ref{c:matrix-presentation} to hold in rational Chow groups
on every power:
\begin{enumerate}[label=(\roman*)]
\item $c_n(T_X)=d\,o$ and $\Delta_*o=o\boxtimes o$;
\item $\bigwedge^{d+1}h(X)=0$ in rational Chow motives;
\item $P_a(h(X),\mu,1,\operatorname{tr},o,\beta^{-1})=0$ for
      $1\leq a\leq s$.
\end{enumerate}
Here $\mu$ is the small-diagonal multiplication,
$\operatorname{tr}:h(X)\to\Q(-n)$ is the trace,
$\beta=\operatorname{tr}\circ\mu$, and
$\beta^{-1}:\Q(-n)\to h(X)\otimes h(X)$ is its inverse pairing.
Every condition is the vanishing of a specified finite linear
combination of correspondence diagrams.  Subject to (i), each can
equivalently be written as a finite list of cycles in the requested
point-and-diagonal rings.  The construction uses finite-dimensional
polynomial algebra, so is a finite criterion, rather than one separate
test for each power of $X$.
\end{theorem}

\begin{proof}
Write $A=H^*(X,\Q)$ as an ordinary vector space; the pure-even
hypothesis makes its tensor symmetry the ordinary one.  Retain the
Tate degree by a one-dimensional auxiliary representation $L$, and
put $G=\operatorname{GL}(A)\times\mathbb G_m$, with the second factor
acting on $L$.  The data
\[
 t=(\mu,1,\operatorname{tr},o,\beta^{-1})
\]
form one rational point of the finite-dimensional rational
$G$-representation
\[
 W=(A\otimes A^\vee\otimes A^\vee)\oplus A
   \oplus(A^\vee\otimes L^n)
   \oplus(A\otimes L^{-n})
   \oplus(A\otimes A\otimes L^{-n}).
\]
In the point $t$, the coordinate denoted by $o$ is the cohomology class
of the degree-one zero-cycle.  Let $I\subset\Q[W]$ be the ideal of
the Zariski closure of $G\cdot t$.  We describe all steps producing
the finite list.

First compute generators of $I$ by elimination: write the orbit map
in matrix entries of $g\in\operatorname{GL}(A)$ and a nonzero scalar
$z$, introduce inverses with equations $u\det(g)=1$ and $vz=1$, and
eliminate $g,u,z,v$ from the graph ideal.  The resulting elimination
ideal is the ideal of the image closure.  Hilbert's basis theorem
makes the resulting list finite.  Enlarge its span to a finite-dimensional
$G$-stable subspace $E\subset I$ generating $I$ as an ideal.
This enlargement is finite because the action on the polynomial ring
is locally finite and preserves the finite degree bound of the chosen
generators.  It can be computed by closing their span under the matrix
units of the Lie algebra and the second torus action; $G$ is connected
in characteristic zero.

The inclusion $E\hookrightarrow\Q[W]$ is equivalently a polynomial
covariant
\[
 p_E\in(E^\vee\otimes\Q[W])^G.
\]
Choose a split $G$-embedding of $E^\vee$ into a finite direct sum $T$
of mixed tensor powers of $A,A^\vee,L,L^\vee$.  Such a split embedding
exists because rational representations of $G$ are semisimple and
are generated by its defining representations and their duals.
For an explicit terminating search, enumerate these finite direct
sums and compute the rational intertwiner spaces.  Enumerate their
rational elements until one has full rank on $E^\vee$, and then
solve the linear intertwiner equations for a retraction.  The
existence of a split embedding proves termination.  The components of the image of
$p_E$ in $T\otimes\Q[W]$ are the covariants we denote by $P_a$.

We recall the two representation-theoretic facts needed here.
Over $\C$, averaging a Hermitian form over
$U(d)\times U(1)$ gives invariant complements.  This compact group
is Zariski dense in $G_\C$, so the complements are $G_\C$-stable.
The intertwiner and splitting equations have rational coefficients;
Gaussian elimination descends their solvability to $\Q$.
Every finite-dimensional rational representation embeds, by its
matrix coefficients, into finitely many copies of $\Q[G]$.
The relevant coefficient functions have bounded polynomial degree
and bounded powers of $\det^{-1}$ and the second torus coordinate.
They therefore occur in mixed tensor powers, using symmetric and
exterior projectors.  Complete reducibility makes the embedding
split.  This proves the existence statements just used.  It also
proves exactness of invariants for the locally finite rational
modules below: place a preimage and its translates in a
finite-dimensional submodule and split the resulting surjection.
The exact references are \citet[Theorem~4.14, pp.~88--89;
Theorem~22.42, pp.~478--479; Remark~12.53(a), p.~248,
and the Reynolds operator on p.~249]{c_milneag2022}.

We explain why these covariants are tensor expressions of the promised
kind.  Polarize a homogeneous polynomial covariant in the finitely
many tensor coordinates of $W$.  Its coefficients are invariants in
mixed tensor powers of $A$ and $A^\vee$.  The first fundamental
theorem says that these are sums of complete contractions, with any
uncontracted indices supplying the target tensor.  The only linear
relations among such contractions are generated by the
$(d+1)$-fold antisymmetrizer and by the scalar loop value $d$.
These are precisely \citet[Theorems~1.2--1.5, pp.~4--5]{c_procesi2021tensor};
polarization and restitution apply because the characteristic is zero.
This use is of the classical tensor theorems, and does not apply a
matrix-only statement to arbitrary tensor variables without that step.
The extra $\mathbb G_m$ factor merely selects the matching Tate
degrees.  Thus each $P_a$ is a finite rational linear combination of
the prescribed diagrams, and every universal identity in their
ordinary $d$-dimensional interpretation follows from the stated
antisymmetrizer and loop relation.

Now let $F$ be any polynomial covariant with target a mixed tensor
space $T'$ that vanishes at $t$.  Equivariance implies that it
vanishes on $G\cdot t$ and its closure.  Hence
\[
 F\in(T'\otimes I)^G.
\]
The multiplication map $E\otimes\Q[W]\twoheadrightarrow I$ is
$G$-equivariant.  Since $G$ is reductive over $\Q$, taking invariants
is exact, and $F$ lifts to
\[
 \widetilde F\in(T'\otimes E\otimes\Q[W])^G.
\]
Contracting its $E$ index against $p_E$ gives $F$ exactly.  Extend
that contraction through the chosen split embedding into $T$.
The first fundamental theorem again writes $\widetilde F$ as a
diagram.  Consequently the tensor ideal of all relations at $t$
is generated by the finite covariants $P_a$, together with the
$(d+1)$-fold antisymmetrizer and the loop relation.  No reductivity
of the stabilizer of $t$ has been assumed or used.

Evaluate this presentation in Chow motives.  The scalar loop value
is $\chi(X)=d$ by the categorical trace formula.  Condition (ii)
annihilates the dimension relation, and condition (iii) annihilates
the finite orbit-ideal generators.  Therefore every cohomological
diagram relation vanishes in Chow, in particular every relation in
the requested rings.  This proves sufficiency.
The additional identities for the auxiliary line,
$\bigwedge^2 L=0$ and $\operatorname{tr}(\id_L)=1$, already hold
for the fixed even invertible Tate object $L=\Q(-1)$; they are
not additional conditions on $(X,o)$.

Conversely suppose all the requested relations hold in Chow.
The cohomological relations
$\Delta\cdot o_1=o_1o_2$ and
$\Delta^2=d\,o_1o_2$ hold there.  The first is the second assertion
of (i); pushing the second to $X$ gives its first assertion.
The alternating projector in (ii) is a rational linear combination
of permutation graphs on $X^{d+1}$, each a product of pairwise
diagonals in $X^{2(d+1)}$.  Its cohomology is zero, so (ii) follows.

Finally the correspondence diagrams used above do not enlarge the
vanishing test after (i).  The Frobenius identities reduce each
connected diagram to a small diagonal with some point insertions
and some handles.  A handle inserts $c_n(T_X)$; a point insertion
inserts $o$.  By (i) both reduce to scalar multiples of $o$, two
such insertions vanish for dimension reasons, and a single insertion
on a small diagonal is the external product of the point classes.
Closed components give their degrees.  Thus every diagram reduces
to the pointed partitions of Definition~\ref{c:def-pointed}.
This is the same graph reduction as Theorem~\ref{c:graph-rule},
now including contractions and using the two verified Chow
identities.  Each $P_a$ has zero cohomology and reduces to a relation
of the requested kind; it consequently vanishes in Chow.  This
proves necessity and the finite criterion.
\end{proof}

\begin{remark}[Scope of the finite criterion]
\label{c:finite-criterion-scope}
\status{Proved scope clarification; the arbitrary-pair classification remains unresolved}
Theorem~\ref{c:finite-even-criterion} applies to arbitrary cup-product
algebras with pure-even cohomology, not only to surfaces.  It gives
a finite algebraic construction of the cycle conditions, without
claiming to decide their rational equivalence by computation.
For general mixed cohomology the analogous direct use of
$\operatorname{GL}_{a|b}$ would require exactness of its invariant
functor, which fails.  A chosen Chow parity splitting would permit
the ordinary reductive group $\operatorname{GL}_a\times
\operatorname{GL}_b$ and the even and odd dimension relations, but
its existence is an additional issue and it need not lie in the
requested tautological correspondence category.  The unrestricted
classification of all $(X,o)$ therefore remains distinct from the
specific K3 obstruction proved below.
\end{remark}


\begin{theorem}[The narrow K3 ring already detects finite-dimensionality]
\label{c:k3-gate}
\status{Proved equivalence; the general vanishing is unresolved}
Let $S/\C$ be a projective K3 surface and let $o$ be its Beauville--Voisin class. The following assertions are equivalent:
\begin{enumerate}[label=(\roman*)]
\item for every $k$, all the cohomological relations among the point and diagonal classes hold in $\CH^*(S^k)_\Q$;
\item with $\pi_2=\Delta_S-o\times S-S\times o$, the cycle
\begin{equation}
\label{c:k3-determinant}
 D_{23}=\sum_{\sigma\in S_{23}}\sgn(\sigma)
          \prod_{a=1}^{23}(\pi_2)_{a,23+\sigma(a)}
 \quad\text{vanishes in }\CH^{46}(S^{46})_\Q;
\end{equation}
\item $\bigwedge^{23}h^2(S)=0$ in rational Chow motives;
\item $h(S)$ is finite-dimensional in the sense of Kimura--O'Sullivan.
\end{enumerate}
Thus the request for every K3 surface is not supplied by the Beauville--Voisin small-diagonal theorem. The first missing universal determinant for this narrower ring occurs on $46$ factors.
\end{theorem}

\begin{proof}
The projectors $o\times S$, $\pi_2$, and $S\times o$ are mutually orthogonal Chow projectors, giving $h(S)=\mathbf1\oplus h^2(S)\oplus\Q(-2)$ (with the contravariant/cohomological convention). The cycle $\pi_2$ belongs to the requested ring. As a correspondence on $S^{23}$, $D_{23}$ is $23!$ times the alternating projector on $h^2(S)^{\otimes23}$. Consequently its vanishing is exactly (iii), because an object in a pseudo-abelian category is zero precisely when its defining idempotent is zero. Its cohomology class is zero since $\dim H^2(S)=22$, so (i) implies (ii). Assertion (iii) makes $h^2(S)$ evenly finite-dimensional, and the other two summands are Tate motives; hence it implies (iv). Conversely, finite-dimensionality makes the even summand have exterior dimension $22$, giving (iii); alternatively the dimension assertion is the finite-dimensionality statement recalled in \citet[\S2.4, p.~506]{c_yin2015}. Finally, \citet[unnumbered main Theorem, p.~503, and \S\S2.5, 3.1--3.6]{c_yin2015} proves that finite-dimensionality is equivalent to injectivity for the larger ring generated by all divisor classes and diagonals. Its restriction gives (i).

We can also see sufficiency directly in the smaller ring. \citet[Theorem 1 and Proposition 3.2]{c_beauvillevoisin2004} gives $c_2(T_S)=24o$ and the small-diagonal relation. These imply every local relation in \eqref{c:pair-local} with $b=22$. The only additional generators supplied by Theorem~\ref{c:pair-ranks} are the permutations and products of \eqref{c:k3-determinant}. This proves that no larger unlisted family of relations is hidden in the equivalence.
\end{proof}

The statement just proved is a reduction, not a proof of (i) for general K3 surfaces. The current-source check includes \citet[Theorem L, explicitly Theorem 9.7 and Corollary 9.8]{c_floccarifu2026}. That result proves abelian Chow motives for K3 surfaces whose transcendental quadratic space embeds in $U_\Q^{\oplus2}\oplus\langle-1\rangle\oplus\langle-b\rangle$, with $b>0$ and $b\equiv3\pmod4$; it supplies new four-dimensional families of generic Picard rank $16$. It does not assert the vanishing of \eqref{c:k3-determinant} for every projective K3 surface. The general classification of all pairs $(X,o)$ requested in C.1(c) also remains a separate open row; the finite presentation above gives its exact cycle conditions but does not turn them into a geometric classification.

\begin{corollary}[Removing the algebraic part of the K3 gate]
\label{c:k3-transcendental-gate}
\status{Proved smaller equivalent vanishing test; general vanishing unresolved}
Let $S/\C$ be a projective K3 surface of Picard rank $\rho$, let
$D_1,\ldots,D_\rho$ be a rational basis of $\operatorname{NS}(S)_\Q$,
and let $B=(\deg D_aD_b)_{a,b}$.  Define
\[
 \pi_{\rm alg}=\sum_{a,b}(B^{-1})_{ab}D_a\times D_b,
 \qquad \pi_{\rm tr}=\pi_2-\pi_{\rm alg},\qquad
 t_2(S)=(S,\pi_{\rm tr},0).
\]
Every assertion of Theorem~\ref{c:k3-gate} is also equivalent to
\[
 \bigwedge^{23-\rho}t_2(S)=0.
\]
Equivalently, the alternating sum of the products
$(\pi_{\rm tr})_{a,,23-\rho+\sigma(a)}$, over
$\sigma\in S_{23-\rho}$, vanishes on $S^{,2(23-\rho)}$.
For Picard rank $1$ this is a test on $44$ factors.  This smaller
correspondence uses divisor classes outside the narrow requested
ring; its factor count therefore does not contradict the unconditional
narrow-ring result through $45$ factors.
\end{corollary}

\begin{proof}
The intersection matrix $B$ is nonsingular by the Hodge index
theorem.  Direct composition of the product correspondences shows
that $\pi_{\rm alg}$ is an idempotent with image $\Q(-1)^{\oplus\rho}$,
and that it is orthogonal to the degree-zero and degree-four
projectors.  Thus
\[
 h^2(S)=\Q(-1)^{\oplus\rho}\oplus t_2(S).
\]
For a direct sum in a rational pseudo-abelian symmetric category,
the exterior-power decomposition gives
\[
 \bigwedge^{23}h^2(S)
 =\bigoplus_{a=0}^{\rho}
   \bigwedge^a\bigl(\Q(-1)^{\oplus\rho}\bigr)
       \otimes\bigwedge^{23-a}t_2(S).
\]
The $a=\rho$ summand is
$\Q(-\rho)\otimes\bigwedge^{23-\rho}t_2(S)$.
Its vanishing follows from that of the whole direct sum; tensoring
with the inverse Tate object proves one implication.
Conversely, every higher exterior power is a direct summand of
the tensor product of the $(23-\rho)$th exterior power with a
further tensor power.  Hence the stated vanishing kills every
summand in the display.  The correspondence form is exactly the
alternating-projector description.  This also checks directly the
transcendental reduction in \citet[\S2.5, pp.~506--507, equation~(7)]{c_yin2015}
while explaining its relation to the smaller ring in the question.
\end{proof}


\begin{corollary}[The unconditional K3 range]
\label{c:k3-stable-range}
\status{Proved}
For every projective complex K3 surface with its Beauville--Voisin class, the requested relations hold in every power $S^k$ with $k\leq45$. More generally they hold in normalized degree $d$ of every $S^k$ whenever
\[
 \min(d,k-d)\leq22,
\]
where normalized degree $d$ means Chow codimension $2d$. In this range no finite-dimensionality assumption is required.
\end{corollary}

\begin{proof}
The Beauville--Voisin identities give the local relations, so every Chow monomial reduces to a linear combination of disjoint matching monomials
$\prod_{i\in I}o_i\prod_{\{a,b\}\in\mathcal M}q_{ab}$.
If the matching has $m$ edges, degree $d$ means $|I|+m=d$, and the number of used factors is $|I|+2m$. Consequently $m\leq\min(d,k-d)$. For $m\leq22$, the orthogonal matching tensors in $(H^2(S))^{\otimes2m}$ are linearly independent: the first determinantal relation requires $23$ edges, as proved in Theorem~\ref{c:pair-ranks}. Monomials with different point and matching supports lie in different K\"unneth multidegrees. Thus their cohomological classes are independent, and no homologically zero Chow combination remains. For $k\leq45$, every matching has at most $22$ edges. This proves both assertions.
\end{proof}

The remaining relation cannot be obtained merely by asserting that all lower-power relations hold. It first appears in the middle normalized degree of $S^{46}$ and is exactly \eqref{c:k3-determinant}. We also tested three possible ways to remove this obstruction. A Kuga--Satake correspondence giving a splitting in homological motives would produce a composite $1+\eta$ in Chow, with $\eta$ homologically trivial; promoting it to a Chow splitting still requires nilpotence or another argument making $1+\eta$ invertible. Specializing to dense loci with finite-dimensional motives does not prove the generic Chow vanishing, because the rational-equivalence locus is only controlled as a countable union of closed loci without a uniform family of equivalences. Finally, the known small-diagonal identities are precisely the local matching reductions and do not supply the first orthogonal determinant. These are failed proof steps, not counterexamples to finite-dimensionality.

The current-source search also checked the official withdrawal notice of the October 2026 paper claiming the Hodge conjecture for all K3 products \citep{c_k3withdrawal2026}. We exclude that withdrawn proof. Even a valid homological Kuga--Satake splitting alone would leave the Chow nilpotence step just identified.

\begin{example}[An explicit genus-three failure]
\label{c:fermat-counterexample}
\status{Proved counterexample}
Let $C/\C$ be the plane quartic $X^4+Y^4=Z^4$, and take $o=(0:1:1)$. Its modified diagonal $M_o$ vanishes in every Weil cohomology, but
\[
 M_o\ne0\quad\hbox{in }\CH^2(C^3)_\Q.
\]
Thus the cohomological relation in Proposition~\ref{c:modified-diagonal} fails in rational Chow groups for this completely specified genus-three pointed curve.
\end{example}

\begin{proof}
The three partial derivatives $4X^3,4Y^3,-4Z^3$ have no common projective zero, so $C$ is smooth. Adjunction for a smooth plane quartic gives $g=(4-1)(4-2)/2=3$. The displayed point lies on $C$. Every product of three positive-degree cohomology classes on a curve is zero; Proposition~\ref{c:modified-diagonal} therefore gives the cohomological vanishing.

Put $J=\operatorname{Jac}(C)$ and $i:C\to J$, $x\mapsto[x-o]$. We use the precise nontorsion assertion of \citet[Theorem 4.1, printed p.~104]{c_bloch1984}: the Ceresa cycle $\Sigma=i_*[C]-[-1]_*i_*[C]$ for this Fermat quartic has infinite order modulo algebraic equivalence. We checked the detected class in its proof: the projected $5$-adic Abel--Jacobi class is nonzero by Theorem 3.4; a hypothetical torsion class would come from invariants of a $\Q_5/\Z_5$ module on which Frobenius at $11$ acts as $-11$, and $-11-1$ is a $5$-adic unit. Restriction--corestriction prevents any finite extension from killing the nontorsion class. The assertion persists over $\C$: finite rational- or algebraic-equivalence data over $\C$ spread over a finitely generated $\overline\Q$-algebra and specialize to a closed point, giving such data over a finite extension, which would contradict the theorem.

We make the passage to $M_o$ directly in rational Chow groups. For this three-dimensional Jacobian, the decomposition theorem and Proposition 3(a) of \citet[pp.~647--649]{c_beauville1986} give
\[
 \CH^2(J)_\Q=\CH^2(J)_{(0)}\oplus\CH^2(J)_{(1)}\oplus\CH^2(J)_{(2)}.
\]
The odd cycle $\Sigma$ lies entirely in grade $1$. Proposition 1(iv) of the same source therefore gives $[m]_*\Sigma=m^3\Sigma$.

For $s:C^3\to J$, $(x,y,z)\mapsto i(x)+i(y)+i(z)$, direct pushforward of the seven terms gives
\[
 s_*M_o=[3]_*i_*[C]-3[2]_*i_*[C]+3i_*[C].
\]
This is also the identity in \citet[Proposition 5.3]{c_grossschoen1995}. Taking its odd part under $[-1]_*$ gives
\[
 \frac{1-[-1]_*}{2}s_*M_o
 =\frac12\bigl([3]_* -3[2]_*+3\bigr)\Sigma.
\]
The identity just proved in the Chow group makes this odd part exactly
\[
 \tfrac12(3^3-3\cdot2^3+3)\Sigma=3\Sigma\ne0.
\]
Hence $s_*M_o$ and therefore $M_o$ are nonzero in rational Chow groups, as claimed.

The polynomial smoothness certificate, the point substitution, the genus value, and the coefficient $3^3-3\cdot2^3+3=6$ are recorded by \texttt{check\_explicit\_bad\_reduction.py}. The nontorsion input is the opened theorem and its proof, not a claim that this finite polynomial certificate computes a Chow group.
\end{proof}

\subsection{C.2: the special bivariant systems}

The relevant category must retain the two maps of manuscript 11:34. Requiring an ordinary algebra homomorphism to preserve both the unit and the Frobenius counit would change that category and would exclude the geometric maps under discussion.

\begin{definition}[A gauged special bivariant system]
\label{c:bivariant-def}
\status{Definition, with the corrected exchange axiom}
Over a characteristic-zero field $L$, a special bivariant system assigns to every finite set $I$ a finite-dimensional graded-commutative algebra $A(I)$, with $A(\varnothing)=L$, a trace $\kappa_I$ of degree $-D|I|$, and a perfect trace pairing. For a map $\alpha:I\to J$ it assigns
\[
 \alpha^*:A(I)\longrightarrow A(J),\qquad
 \alpha_*:A(J)\longrightarrow A(I).
\]
The first is a unital graded algebra homomorphism; the second has degree $D(|I|-|J|)$. Pullbacks compose covariantly in maps of indexing sets, pushforwards contravariantly, and
\[
 \kappa_I(\alpha_*x\,a)=\kappa_J(x\,\alpha^*a),\qquad
 \alpha_*(x\alpha^*a)=\alpha_*x\,a.
\]
For disjoint unions the exterior-product maps $A(I)\otimes A(J)\to A(I\amalg J)$ are isomorphisms compatible with associativity, units, traces, and the Koszul symmetry. These maps are compatible with pullback and pushforward. The exchange axiom is base change for product squares (a projection pulled back by a map on the other factor); arbitrary self-intersection squares are not included. The gauge is the chosen trace, not a quotient up to multiplication by a scalar.

A morphism in the bivariant category from $A$ to $B$ is a pair $(f^*,f_*)$, with $f^*:B\to A$ a unital graded algebra map and $f_*:A\to B$ its trace adjoint. On powers these maps are tensor powers and satisfy the corresponding pullback and pushforward compatibilities. The adjoint need not preserve the unit.
\end{definition}

\begin{theorem}[C.2: equivalence and the exchange axiom]
\label{c:bivariant-equivalence}
\status{Proved, with the category specified above}
Special gauged bivariant systems are equivalent to finite-dimensional graded-commutative Frobenius algebras of even top degree, with bivariant adjoint morphisms. The functor taking a system to its one-factor algebra and diagonal coevaluation is fully faithful. For a smooth projective connected variety, $I\mapsto H^*(X^I)$ satisfies these axioms. Correspondence composition defined by these axioms is associative and has the diagonal as identity.
\end{theorem}

\begin{proof}
From a Frobenius algebra $(A,\mu,1,\kappa)$, set $A(I)=A^{\otimes I}$. For $\alpha:I\to J$, let $\alpha^*$ multiply the inputs along each fibre $\alpha^{-1}(j)$, inserting $1$ for an empty fibre, with the Koszul permutation signs. Associativity and commutativity of $\mu$ give functoriality. Define $\alpha_*$ as its unique trace adjoint. Perfectness gives existence and uniqueness; taking adjoints proves pushforward functoriality. Pairing both sides against an arbitrary element proves the projection formula. Product-square base change follows by separating the tensor factors: contraction in the projected factor commutes with the map on the other factor. This also verifies all unit and exterior-product compatibilities.

Conversely, specialness identifies $A(I)$ with $A(\{1\})^{\otimes I}$. A fold of two indices is the product, the empty fibre gives its unit, and the trace gives a perfect pairing. Every map of finite sets is a composite of folds, injections, and bijections. Its pullback is therefore forced by this algebra; its pushforward is forced by adjointness. The coevaluation $\Delta$ is the unique tensor whose contraction is the identity. This proves reconstruction and full faithfulness. A bivariant pair on one factor extends to tensor powers, and the required compatibilities follow respectively from multiplicativity and its adjoint. We do not impose the generally false extra identity $(f^*\otimes f^*)\Delta_B=\Delta_A$.

For Weil cohomology, K\"unneth, Poincar\'e duality, and the projection formula give exactly these maps. A map $\alpha:I\to J$ induces the geometric map $X^J\to X^I$, so all variances agree with the displayed formulas. Product base change is again K\"unneth and trace in the omitted factor. In contrast, pulling a diagonal back through itself produces its Euler self-intersection, so unrestricted Cartesian base change without an excess term is false. Refined Gysin base change may be included on transverse squares, with the transversality hypothesis retained.

Finally, composition of kernels is tensor contraction in the middle factor. Two successive contractions on different factors commute by associativity of tensor product and the Koszul rule. The resulting two contractions are exactly the two parenthesizations of a triple composition. Contraction with the coevaluation is the identity by duality. This proves the associativity requested in 11:48. It also supplies the precise exchange axiom behind ``axiome cart\'esien oubli\'e !'' (11:46), and answers ``il convient de pr\'eciser les axiomes'' (11:42) on the special objects actually used here.
\end{proof}

\begin{remark}[The R7 hypotheses used in this Part]
\label{c:r7-audit}
\status{Used and rechecked, with its internal-Hom counterexample confirmed}
We checked the proof in the supplied reading, \S\S2.1--2.3, against 12:150--158. The conditional construction there is not a construction of a faithful realization on all Chow motives: it first quotients by the chosen homological equivalence, assumes its equality with numerical equivalence and the stated positivity, and then takes the semisimple envelope. Our Chow assertions above do not use this conditional passage. In particular it cannot be used to kill the homologically zero K3 determinant.

The sign correction is forced by the cross-product formula: interchanging a class of degree $a$ with one of degree $b$ multiplies it by $(-1)^{ab}$. Thus the original realization is symmetric into graded vector spaces with that commutativity constraint. This is the convention used in our Frobenius equivalence. A level filtration by thick subcategories yields a direct-sum level decomposition only when the ambient category is semisimple; then we group the simple constituents according to their least allowed twist. Thickness alone gives no splitting of extensions. The geometric level definition of C.4 therefore does not assert a canonical splitting or identify itself with coniveau.

The internal-Hom error of 12:156 already occurs in the semisimple effective Tate category with objects finite sums of $\Q(-a)$, $a\geq0$. For every effective $Z$,
\[
 \Hom(Z\otimes\Q(-1),\Q)=0.
\]
Hence the effective internal Hom from $\Q(-1)$ to $\Q$ is the zero object. After adjoining all Tate twists the corresponding internal Hom is $\Q(1)$, which is nonzero. Thus effective internal Homs need not be preserved by Tate inversion. In a rigid tensor category the valid identity is
\[
 \underline{\Hom}(M,N(-i))
 =M^\vee\otimes N(-i)
 =M^*\otimes N,
 \qquad M^*=M^\vee(-i).
\]
Evaluation and coevaluation give these isomorphisms and their inverse maps. Merely postulating the effective object $M^*$ without the duality identities does not prove them. These checks confirm precisely the restrictions of R7 needed here.
\end{remark}

\subsection{C.3: the algebraic part of the Lie object}

We separate the Lie algebra of an algebraic tannakian group from the analytic Lie algebra of a compact $\ell$-adic image. They are different objects. In a neutral tannakian category over $\Q$, the algebraic part of an adjoint Lie object means its maximal trivial subobject, computed as $\Hom(\mathbf1,\Lie G)$. The same definition is made after neutralizing a gerbe, with its descent action retained.

\begin{proposition}[The invariant-centre identity]
\label{c:lie-centre}
\status{Proved}
For a pro-reductive group $G$ over a characteristic-zero field,
\[
 \Hom_G(\mathbf1,\Lie G)=(\Lie G)^G=\Lie Z(G).
\]
If $G$ is connected, the map from its centre to its abelianization induces an isomorphism on Lie algebras. If $G^\circ$ has protorus abelianization $S$ and $\Gamma=G/G^\circ$, then
\[
 (\Lie G)^G=(\Lie S)^\Gamma.
\]
The invariant condition under $\Gamma$ cannot be discarded.
\end{proposition}

\begin{proof}
An infinitesimal element fixed by every conjugation is exactly an infinitesimal central element: the equations defining these two closed subgroup functors give the same kernel after differentiation. For a connected reductive group, $Z(G)^\circ\times G^{\mathrm{der}}\to G$ is a central isogeny. Its differential is an isomorphism, and the semisimple derived Lie algebra has no adjoint invariants. This proves the connected statement. The connected component acts trivially on its centre; the remaining action is precisely that of the component group. Passing to the inverse limit of finite-dimensional reductive quotients gives the pro-reductive assertions, since both Lie algebras and fixed-vector conditions are computed compatibly in those quotients.
\end{proof}

\begin{theorem}[Lead C3 with its arithmetic action]
\label{c:serre-centre}
\status{Proved under the stated Serre-quotient hypothesis}
Let $K$ be finitely generated over $\Q$, embedded in $\C$, and let $K_0$ be the algebraic closure of $\Q$ in $K$. Suppose that the geometric connected motivic group is pro-reductive, that its maximal protorus quotient is the full Serre protorus $S$, and that the action of $\Gal(\overline K/K)$ on this quotient is the usual action through $\Gal(\overline\Q/K_0)$. If $E$ is the largest CM subfield of $K_0$, then
\[
 \dim_\Q (\Lie G)^G=
 \begin{cases}1+[E:\Q]/2,&E\text{ exists},\\1,&K_0\text{ has no CM subfield}.
 \end{cases}
\]
These hypotheses hold for the full category of motivated motives of abelian type, including Artin motives. The conclusion concerns the full arithmetic adjoint invariants; the centre of the geometric connected group alone is generally an infinite-dimensional pro-vector space.
\end{theorem}

\begin{proof}
We first prove the character calculation rather than infer it from the dimension of a geometric torus. Write $\Q^{\mathrm{CM}}$ for the union of finite Galois CM extensions of $\Q$. Complex conjugation is central in its Galois group. The rational character space of the Serre protorus is the union of the spaces of locally constant rational-valued functions on the finite Galois quotients satisfying
\[
 f(\sigma)+f(c\sigma)=w(f),\qquad w(f)\text{ constant}.
\]
It decomposes as the one-dimensional constant space plus the space of $c$-odd functions. Taking coinvariants under $\Gamma_{K_0}$ is exact at each finite level, because averaging over the finite image identifies coinvariants with invariants in characteristic zero. At a finite Galois CM field $L$, invariant functions are functions on the embeddings of $L\cap K_0$. If this intersection is totally real, conjugation fixes every embedding and no odd function remains. If it is CM, conjugation pairs its embeddings without fixed points and the odd space has dimension $[L\cap K_0:\Q]/2$.

Every subfield of a Galois CM field is either totally real or CM: central conjugation descends to it, and the fixed field is totally real. The CM subfields of the finite number field $K_0$ have a compositum which is again CM, hence a largest one $E$ whenever any exists. Taking $L$ containing its normal closure makes the preceding dimensions stabilize at the asserted value. Thus the coinvariant character space is finite-dimensional of that dimension. Since $\Lie S=\Hom_\Q(X^*(S)\otimes\Q,\Q)$ as a pro-vector space, its fixed part is dual to those coinvariants. The preceding proposition proves the formula.

Restriction $\Gamma_K\to\Gamma_{K_0}$ is surjective: $K/K_0$ is regular, hence linearly disjoint from every finite extension of $K_0$, and the compatible extension property for finite Galois extensions gives surjectivity on inverse limits. Therefore no extra constants enter the invariant calculation.

For motivated motives of abelian type, \citet[Theorem 0.6.2 and the remark following its proof in \S6]{c_andre1996} proves that all Hodge classes on abelian varieties are motivated, including the absolute-Hodge classes over an arbitrary characteristic-zero base. Consequently the geometric Hodge realization is fully faithful on this category; compare \citet[Theorem 8.25 and Proposition 8.23]{c_delignehodge1982}. Its maximal subcategory with torus group is the category of CM Hodge structures, generated by CM abelian varieties. The Serre quotient and its arithmetic action are identified by \citet[unnumbered Theorem, p.~262, and Lemmas 1--2]{c_delignetaniyama1982}. A CM object defined over a finite extension is brought into the category over $K$ by restriction of scalars and then recovered after that extension as a direct summand; thus this assertion uses the full abelian-type category, not merely the CM varieties already having all endomorphisms over $K$. These facts establish precisely the hypotheses used in our character calculation. They also explain why ``no CM subfield'' does not mean that the geometric category contains no CM objects.
\end{proof}

\begin{example}[The finite-field analytic/algebraic distinction]
\label{c:finite-field-lie}
\status{Proved counterexample to the one-dimensional algebraic reading}
Let $E/\F_5$ have equation $y^2=x^3-x$. It is smooth since its discriminant $64$ is nonzero modulo $5$. It has eight rational points, so geometric Frobenius has polynomial $T^2+2T+5$ on $H^1$ and eigenvalues $\pi=-1+2i$, $\bar\pi=-1-2i$. Its algebraic Frobenius torus is
\[
 T=\operatorname{Res}_{\Q(i)/\Q}\Gm,
 \qquad\dim T=2.
\]
Indeed the two primes of $\Q(i)$ above $5$ give valuations $(1,0)$ and $(0,1)$ on $(\pi,\bar\pi)$; hence $\pi^a\bar\pi^b=1$ implies $a=b=0$. There is consequently no character cutting out a proper algebraic subgroup containing Frobenius. The adjoint action of $T$ is trivial, so its whole Lie algebra is algebraic, but it is two-dimensional.

On the other hand, for $\ell=3$ a suitable positive power of $\pi$ belongs to $1+3\Z_3[i]$. The analytic closure of the cyclic Frobenius subgroup has Lie algebra $\Q_3\log(\pi^N)$, of dimension one. The logarithm is nonzero because $\pi$ is not a root of unity. This is the distinction made by Serre in his letter of 8 September 1964, printed p.~202 of \citet{c_grothendieckserre}. Thus ``tout'' on 12:59 is valid for the adjoint Lie algebra of a commutative algebraic motivic band, as in the finite-field category under the Tate hypotheses of Part A; the accompanying identification with a single $\Q(0)$ does not describe that algebraic Lie object. Treating an analytic one-dimensional image as the motivic Lie algebra is an error, not an unconditional proof of a finite-field motivic assertion.
\end{example}

\subsection{C.4: weights, levels, and the overstrengthened bound}

\begin{definition}[The level convention of 12:111--115]
\label{c:level-def}
\status{Definition, checked against 12:111--115}
We fix $\ell$ invertible on the base and use $\Q_\ell(1)$ of weight $-2$. In the dimension-normalized convention, a special lisse sheaf of weight $r\geq0$ is a lisse subquotient of $R^rg_*\Q_\ell$ for a smooth projective morphism $g$ of relative dimension $r$. This includes the inserted dimension clause on 12:111; the transcription marks its following qualifier as uncertain. ``Strictly special'' means a direct summand, as clarified in 12:113. We explicitly call the version without a restriction on the relative dimension the broader convention. Over a field the two conventions agree, by the cut-and-pairing argument in Lemma~\ref{c:proper-level}; Theorem~\ref{c:general-level-bound} first uses the dimension-normalized convention, and Proposition~\ref{c:broader-coefficients} removes that clause for the general polygon and for the requested margin bound. Effective admissibility means that, after a finite stratification and the finite extensions allowed in 12:111, there is a finite filtration by extensions by zero of such sheaves. A sheaf has levels at most $\nu$ if these pieces may be chosen as $V(-a)$ with $a\geq0$, where $V$ is special of weight $r\leq\nu$. The simultaneous condition of weights in $\Phi$ requires the same pieces to satisfy $r+2a\in\Phi$. This is the geometric generation definition in 12:112; it is not simply a definition in terms of the eigenvalue weights, and it is not automatically the coniveau filtration.
\end{definition}

\begin{proposition}[C.4: the general $n-1$ bound is false]
\label{c:level-counterexample}
\status{Refuted, with the precise manuscript correction}
In the generality requested in the prompt, the bound
\[
 \operatorname{level}(R^if_*F)_{\bar y}
 \leq\min\{i,d_y(f)-1\}+j
\]
is false. Take $f:E\to\Spec\F_5$ for the elliptic curve in Example~\ref{c:finite-field-lie}, $F=\Q_\ell$, $j=0$, $i=1$, and any $\ell\ne5$. Then $d_y(f)=1$ and the right side is zero, whereas $H^1(E_{\overline\F_5},\Q_\ell)$ has level exactly one.
\end{proposition}

\begin{proof}
The sheaf $F$ is strictly special of weight zero and $E$ is regular, smooth and projective. Its nonzero $H^1$ is special of weight one, so its level is at most one. A level-zero object, by Definition~\ref{c:level-def}, has successive constituents $V(-a)$ with $V$ of weight zero and $a\geq0$; all its weights are even. The characteristic polynomial $T^2+2T+5$ has two roots of absolute value $\sqrt5$, so this $H^1$ is pure of odd weight one and cannot have level zero. The program records all eight points and the polynomial, so the counterexample does not depend on a catalogue lookup.

The supplied transcription of the final general-case box of 12:120 reads $\min(i,n)+j$, not $\min(i,n-1)+j$. The margin asks ``peut on remplacer par $n-1$ si $f$ de dim $\leq n-1$ ?'' The elliptic structural morphism has relative dimension $n$, so it does not refute that additional margin hypothesis. We therefore refute the strengthened prompt, not an assertion that the manuscript actually made without the condition.
\end{proof}

\begin{proposition}[Sharp elementary cases]
\label{c:level-special-cases}
\status{Proved in the stated cases}
If $f:X\to\Spec k$ is smooth projective of dimension $n$ and $F=\Q_\ell$, then $H^i(X)$ is pure of weight $i$ over a finite field and has level at most $\min(i,2n-i)$. Over a smooth pair $(Y,D)$ with a strict normal-crossing divisor and open immersion $j:Y\setminus D\hookrightarrow Y$, if $r$ branches meet at a geometric point $\bar y$, then
\[
 (R^ij_*\Q_\ell)_{\bar y}\simeq
 \bigwedge^i\bigl(\Q_\ell(-1)^r\bigr).
\]
It vanishes for $i>r$ and, when nonzero, has weight $2i$ and level zero. If a strictly special lisse sheaf extends across this pair, the projection formula tensors this description with its stalk.
\end{proposition}

\begin{proof}
For $i\leq n$, the cohomology itself is special of weight $i$, hence has level at most $i$. For $i>n$, hard Lefschetz identifies it with $H^{2n-i}(X)(-(i-n))$, giving the second bound. Products of ordinary elliptic curves over finite fields attain these geometric level bounds; the Frobenius-slope proof is given in Proposition~\ref{c:level-optimality}. The normal-crossing formula is local: in strict henselian coordinates each missing coordinate hyperplane supplies its Kummer class in $\Q_\ell(-1)$, and the products of distinct Kummer classes form the exterior algebra. Purity and the localization sequence prove this by induction on the number of coordinate divisors. Frobenius acts with weight two on each Kummer class. For an extending lisse sheaf, the stalkwise projection formula applies and gives the last assertion.
\end{proof}

\begin{definition}[The local dimension on 12:120]
\label{c:local-dimension}
\status{Definition, with the punctured-trait correction verified}
For a finite-type morphism $f:X\to Y$ of universally catenary noetherian schemes we use
\[
 d_y(f)=\sup_{\substack{x\in X\\ f(x)\text{ generalizes }y}}
 \left\{\operatorname{trdeg}_{k(f(x))}k(x)
       +\dim\overline{\{f(x)\}}^{\,\Spec\mathcal O_{Y,y}}\right\}.
\]
The closure in the second summand is taken in the local base. This is the dimension function, normalized to be zero at $y$, pulled back to $X$. It is the formula visible in the supplied transcription, whose notation for the local spectrum is uncertain. It measures local source dimension, not maximum fibre dimension: for $\Spec K\hookrightarrow\Spec R$ with $R$ a discrete valuation ring and $y$ its closed point it equals $1$, although all nonempty fibres have dimension zero. After an equidimensional localization of the source and a compactification, its maximum is the local dimension of that compactification. Regular alterations preserve this maximum. The dimension-function convention and the corresponding local cohomological interpretation are made precise in \citet[\S1.1.1 and Corollary 1.2.2]{c_pillonistroh2014}.
\end{definition}

\begin{lemma}[Proper cohomology with simultaneous geometric bounds]
\label{c:proper-level}
\status{Proved}
Let $Z$ be proper of dimension at most $d$ over a field $k$, and let $a\geq0$. After the finite field extensions allowed in the admissibility definition, $H^a(Z_{\bar k},\Q_\ell)$ has a finite filtration by subquotients of special objects $V_r(-t)$ satisfying
\[
 t\geq\max(0,a-d),\qquad r\geq0,\qquad r+2t\leq a.
\]
In particular its weights are at most $a$ and its geometric levels are at most $\min(a,2d-a)$. A negative upper level bound means the object is zero. These are geometric generation assertions, not deductions from a Hodge polygon.
\end{lemma}

\begin{proof}
For a smooth projective variety of dimension $d$, hard Lefschetz identifies $H^a$ with $H^{2d-a}(-(a-d))$ when $a>d$; for $a\leq d$ use $H^a$ itself. The input is \citet[Theorem 4.1.1]{c_deligneweil1980}. If a special object of degree $r$ is required to be presented by a variety of dimension exactly $r$, take a smooth complete intersection of dimension $r$ in the smooth projective variety. Weak Lefschetz injects $H^r$ into its middle cohomology. The restriction pairing is nondegenerate by hard Lefschetz on the original variety, so this injection even has a Galois-equivariant left inverse. A finite extension of $k$ suffices to choose the hyperplanes. Thus both conventions of specialness give the same objects over a field.

For an integral proper $Z$, \citet[Theorem 7.3 and Corollary 7.4]{c_dejong1996}, applied over an algebraic closure and then descended to a finite extension, give a smooth projective $W$, a finite group $G$ acting on it, and a proper map $V=W/G\to Z$ which is a universal homeomorphism over a dense open $U\subset Z$. The construction and the finite group action descend because they involve finitely many algebraic coefficients. Passing through finite purely inseparable extensions does not change the \(\ell\)-adic topos. Let $S=Z\setminus U$ and $E=V\times_ZS$. Both have dimension at most $d-1$. Proper base change identifies
\[
 H^a(V,\Q_\ell)=H^a(W,\Q_\ell)^G:
\]
the finite quotient has no higher direct images of the constant sheaf, and the invariant functions on every geometric orbit are the constants. Averaging is exact over $\Q_\ell$, even if $\ell$ divides $|G|$.

The localization triangles for $U\subset Z$ and its inverse image, together with proper base change, give the abstract blowup sequence
\[
 H^{a-1}(E)\longrightarrow H^a(Z)
 \longrightarrow H^a(V)\oplus H^a(S).
\]
The induction hypothesis for $E$ gives $t\geq\max(0,(a-1)-(d-1))$ and weight at most $a-1$; for $S$ it gives an even stronger lower bound on $t$. The smooth-projective assertion for $W$ handles $V$. Subquotients and extensions preserve all three inequalities. Reducible proper schemes follow by the same localization argument for their components; nilpotents do not change the \(\ell\)-adic topos. This completes the dimension induction.

Finally finite-extension descent does not enlarge the bounds. For a finite separable extension $k'/k$, a representation is a direct summand of the induction of its restriction, by the trace divided by $[k':k]$. A smooth projective $k'$-scheme viewed as a $k$-scheme is still smooth projective of the same dimension, and realizes that induced cohomology. Purely inseparable extensions have already been handled. This justifies the descent used above.
\end{proof}

\begin{theorem}[Corrected bounds and the margin question]
\label{c:general-level-bound}
\status{Proved in the stated arithmetic finite-type setting}
Let $f:X\to Y$ be separated of finite type, with $X$ regular and $X,Y$ of finite type over $\Z[1/\ell]$ or over a field of characteristic different from $\ell$. Fix $y\in Y$ and put $n=d_y(f)$. Let $F$ be strictly special of weight $j$, in the dimension-normalized sense of the insertion on 12:111: it is a direct summand of $R^jg_*\Q_\ell$ for a smooth projective $g$ of relative dimension $j$. Set $\epsilon=1$ if no component of local dimension $n$ maps to $y$, and set $\epsilon=0$ otherwise. Put
\[
 L=\max(0,i-n+\epsilon),\qquad U=\min(i,n)+j.
\]
The stalk $(R^if_*F)_{\bar y}$ is effectively admissible, with a filtration into pieces of weights $w$ and levels bounded simultaneously by
\begin{equation}
\label{c:weight-level-polygon}
 2L\leq w\leq2U,
 \qquad \operatorname{level}\leq\min\{w-2L,\,2U-w\}.
\end{equation}
Consequently
\[
 \operatorname{weight}\leq2\min(i,n)+2j,
 \qquad
 \operatorname{level}\leq
 \begin{cases}
 \min(i,2n-i)+j,&\epsilon=0,\\
 \min(i,2n-1-i)+j,&\epsilon=1.
 \end{cases}
\]
For an open immersion at a boundary point, $\epsilon=1$. The condition that $f$ have relative dimension at most $n-1$ near $y$ also implies $\epsilon=1$. Hence the margin's replacement by $n-1$ is valid with its stated extra condition, since $\min(i,2n-1-i)\leq\min(i,n-1)$. Without that condition the corrected general bound is at most $\min(i,n)+j$, and Proposition~\ref{c:level-counterexample} prevents replacing $n$ by $n-1$ universally.
\end{theorem}

\begin{proof}
We first take $F=\Q_\ell$ and $j=0$. Work locally on $Y$ around $y$, compactify $X$ over $Y$, and retain only the closure of $X$. Applied to that compactification and its boundary, the alteration theorems of \citet[Theorem 4.1 for varieties over fields; Theorem 8.2 for arithmetic schemes]{c_dejong1996} give a proper regular alteration $h:\widetilde X\to Y$ with a strict normal-crossing boundary $D$, and an open $v:X'\hookrightarrow\widetilde X$ above $X$. The induced $\pi:X'\to X$ is proper and generically finite. The construction can be made component by component. Localization to $\Spec\mathcal O_{Y,y}$ preserves regularity, boundary codimensions, and the local dimension function.

There are maps
\[
 \Q_{\ell,X}\longrightarrow R\pi_*\Q_{\ell,X'}
 \xrightarrow{\operatorname{Tr}_\pi}\Q_{\ell,X}
\]
whose composite is multiplication by the generic degree on each component. Here the trace is not an unjustified trace for a nonfinite map: a projective generically finite map between regular schemes is a smoothable local complete-intersection morphism of virtual dimension zero. Absolute purity identifies its Gysin map with $\Q_\ell\to\pi^!\Q_\ell$, and proper adjunction gives the displayed trace. The degree computation is checked on the generic finite-flat locus, since an endomorphism of the constant sheaf on a connected scheme is a scalar. The precise construction and calculation are \citet[Corollary 3.1.2, Lemmas 3.5.8--3.5.10, and Proposition 2.5.13]{c_riou2014}; the equality also covers inseparable generic degree by the finite-flat degree formula. Dividing by the degree proves that $Rf_*\Q_\ell$ is a direct summand of $Rh_*Rv_*\Q_\ell$.

Let $D^{(q)}$ denote the disjoint union of the intersections of $q$ distinct boundary components, and put $D^{(0)}=\widetilde X$. Absolute purity gives
\[
 R^qv_*\Q_\ell=\Q_{\ell,D^{(q)}}(-q).
\]
This is \citet[Corollary 3.1.4]{c_riou2014}, equivalently \citet[Theorem 2.1.2]{c_pillonistroh2014}; it remains true for regular arithmetic schemes. The stalkwise Leray sequence and proper base change now give
\[
 E_2^{a,q}=H^a(F_q,\Q_\ell)(-q)
 \Longrightarrow (R^{a+q}(f\pi)_*\Q_\ell)_{\bar y},
 \qquad F_q=D^{(q)}\times_Y\bar y.
\]
Every $F_q$ is proper. By the definition of $n$ and regular codimension, $\dim F_q\leq n-q$ for $q\geq1$, while $\dim F_0\leq n-\epsilon$. This step does not assume that intersecting a special fibre with a vertical divisor lowers its dimension: the inequality already holds for the entire local scheme $D^{(q)}$, of codimension $q$ in a regular local source of dimension at most $n$. The extra drop for $F_0$ follows because a proper closed fibre containing a component of maximal local dimension would make that component map to $y$. Alteration does not change the image of its generic point.

Write $a=i-q$. Lemma~\ref{c:proper-level} says that each constituent of $E_2^{a,q}$ is $V_r(-(t+q))$, with
\[
 r+2t\leq i-q,
 \qquad t\geq\max(0,i-q-\dim F_q).
\]
For $q=0$ this gives $t\geq L$. For $q\geq1$ it gives $t+q\geq q+\max(0,i-n)\geq L$. Its total weight $w=r+2(t+q)$ satisfies $w\leq i+q$, and nonvanishing requires $q\leq\min(i,2n-i)$ (as well as $q\leq n$). It follows that $w\leq2\min(i,n)$ and $r\leq w-2L$. The other inequality $r\leq2\min(i,n)-w$ follows as follows. If $i\leq n$, use $w\leq i+q$ to obtain $r=w-2(t+q)\leq2i-w$. If $i>n$, use $t\geq i-n$ for $q\geq1$, and combine $q\geq w-i$ with the same formula, obtaining $r\leq2n-w$; the $q=0$ term obeys it because $w\leq i$ and $t\geq i-n+\epsilon$. Thus every term lies in \eqref{c:weight-level-polygon}. A spectral sequence gives a filtration by subquotients of its terms; the inequalities survive extensions and direct summands. No degeneration of this Leray sequence is needed.

For the coefficient sheaf, relative hard Lefschetz and \citet[Theorem 1.5 and Proposition 2.4]{c_deligne1968} give a derived splitting
\[
 Rg_*\Q_\ell\simeq\bigoplus_bR^bg_*\Q_\ell[-b].
\]
Accordingly $R^if_*F$ is a direct summand of $R^{i+j}(fg)_*\Q_\ell$. The source of $g$ is regular, its local dimension is $n+j$, and the value of $\epsilon$ is unchanged. Applying the constant-coefficient result with $(i,n)$ replaced by $(i+j,n+j)$ gives exactly \eqref{c:weight-level-polygon} and the displayed bounds. This proves the corrected statement, including the margin condition.
\end{proof}

\begin{proposition}[Removing the dimension clause from the coefficient]
\label{c:broader-coefficients}
\status{Proved}
Keep the finite-type hypotheses of Theorem~\ref{c:general-level-bound}, but suppose only that $F$ is a direct summand of $R^jg_*\Q_\ell$ for a smooth projective $g:Z\to X$ of arbitrary relative dimension. Put $n=d_y(f)$ and $L_0=\max(0,i-n)$, $U=\min(i,n)+j$. Then the simultaneous bounds
\[
 2L_0\leq w\leq2U,\qquad
 \operatorname{level}\leq\min(w-2L_0,2U-w)
\]
still hold. In particular the general level bound is $\min(i,2n-i)+j$.
If no component of local dimension $n$ maps to $y$, there is also the bound
\[
 \operatorname{level}\leq\min(i,n-1)+j.
\]
Consequently the requested weight bound and the manuscript's margin bound hold in this broader convention as well. The stronger margin polygon with $L_0$ replaced by $\max(0,i-n+1)$ remains covered by Theorem~\ref{c:general-level-bound} in its dimension-normalized convention.
\end{proposition}

\begin{proof}
We first explain the relative dimension reduction on a dense open set. We work componentwise, so the relative dimension $r$ of $g$ is constant. If $r\leq j$, replace $Z$ by $Z\times\mathbb P^{j-r}$; K\"unneth realizes $R^jg_*\Q_\ell$ as a direct summand of the degree-$j$ cohomology of this dimension-$j$ family. If $r>j$, choose over the generic point of $X$ a smooth complete intersection of relative dimension $j$ in a projective embedding of $Z$. It spreads to a smooth projective family $h:W\to U$ over a dense open $U\subset X$. The generic field is infinite unless the source component is zero-dimensional over a finite field; in that case a finite field extension permits Bertini and restriction/corestriction gives the same direct-summand conclusion over the original field.

Weak Lefschetz gives an injection $R^jg_*\Q_\ell|_U\to R^jh_*\Q_\ell$. The pairing induced on the source by middle-dimensional Poincar\'e duality on $W/U$ is cup product with the $(r-j)$th power of the hyperplane class followed by trace. Hard Lefschetz makes it nondegenerate. Taking its adjoint and inverse gives a morphism of lisse sheaves which retracts the injection. Thus $F|_U$ has the dimension-$j$ presentation required by the preceding theorem.

We prove the general polygon by induction on the inherited local dimension. Choose the dense open $U$ just constructed and write $D=X\setminus U$. Since $X$ is regular and excellent, $D$ has a finite stratification by regular locally closed schemes $Z_\alpha$, each of pure codimension $c_\alpha\geq1$ in its source component, with a filtration by closed supports. The dimension formula gives
\[
 d_y(f|_{Z_\alpha})\leq n-c_\alpha.
\]
For nested closed supports $D'\subset D''$, the localization triangle for $R\Gamma_{D''}F$ gives the quotient term obtained by pushing the supported complex on $D''\setminus D'$ from $X\setminus D'$ to $X$. Absolute purity therefore gives, for our regular strata, successive terms
\[
 R(i_\alpha)_*(F|_{Z_\alpha})(-c_\alpha)[-2c_\alpha],
 \qquad i_\alpha:Z_\alpha\longrightarrow X.
\]
Here $i_\alpha$ is the locally closed immersion, and the pushforward is derived ordinary pushforward. Purity applies after the open restriction in which that immersion is closed; its exact form is \citet[Corollary 3.1.2]{c_riou2014}. Composing with $Rf_*$, then using the induction hypothesis on $Z_\alpha$, bounds the supported terms.

Explicitly write $c=c_\alpha$, $n'=d_y(f|_{Z_\alpha})$ and $m=i-2c$. If $m<0$ there is no term. Otherwise the induction hypothesis supplies pieces $V_a(-t_0)$ with
\[
 t_0\geq\max(0,m-n'),\qquad a+t_0\leq\min(m,n')+j.
\]
The additional twist replaces $t_0$ by $t=t_0+c$. Since $n'\leq n-c$,
\[
 t\geq\max(c,i-n)\geq L_0,
 \qquad a+t\leq\min(i-c,n)+j\leq U.
\]
These are exactly the asserted simultaneous polygon inequalities. The open term $R(f|_U)_*(F|_U)$ satisfies them by Theorem~\ref{c:general-level-bound}. The localization triangle $R\Gamma_DF\to F\to Rv_*F|_U$, with $v:U\hookrightarrow X$, now proves them for $Rf_*F$: in each cohomological degree the middle group is an extension of a quotient of the first group by a subobject of the third. Finite extensions, subquotients and filtrations preserve the two inequalities. The zero-dimensional starting case was the field reduction above. This proves the general polygon.

For the extra margin bound we use a different descent that does not introduce vertical closed strata. The smooth-complete-intersection parameter space is a nonempty open in a product of relative projective spaces and is smooth and surjective over $X$. A smooth surjection has sections after an \emph{\'etale} covering. Quasi-compactness gives a finite-type surjective \'etale map $v:V\to X$ over which the coefficient has a dimension-$j$ presentation. Denote the terms of its \v Cech nerve by $v_a:V_a\to X$. They are \'etale over $X$, hence regular, and their inherited local dimensions are at most $n$. When that dimension is $n$, no component maps to $y$; when it is smaller, it is at most $n-1$. The normalized theorem consequently gives, for every $a,b$,
\[
 \operatorname{level}\bigl(R^b(fv_a)_*v_a^*F\bigr)_{\bar y}
       \leq\min(b,n-1)+j.
\]
The \'etale \v Cech descent spectral sequence is
\[
 E_1^{a,b}=\bigl(R^b(fv_a)_*v_a^*F\bigr)_{\bar y}
       \ \Longrightarrow\ (R^{a+b}f_*F)_{\bar y}.
\]
For torsion sheaves this is \citet[Proposition 96.21.1, Tag 072G; see also Theorem 59.19.2, Tag 03OW]{c_stacksdescent}. For a lisse $\Q_\ell$-sheaf choose a stable lattice, apply it modulo $\ell^s$, and pass to the inverse limit and then invert $\ell$. Constructibility makes the finite-level stalks finite, so their inverse systems are Mittag--Leffler. Each total degree has only finitely many relevant first-quadrant terms, which justifies this passage and convergence. Since $b\leq a+b=i$, every contributing piece has level at most $\min(i,n-1)+j$. This proves the margin assertion. No properness of the \'etale cover was needed, and no stronger reflected margin inequality is inferred from this spectral sequence.
\end{proof}

\begin{proposition}[Optimality and the limit of the word sharp]
\label{c:level-optimality}
\status{Proved optimality in the stated senses}
For constant coefficients the preceding upper weight and level bounds are attained for every nonnegative degree in the indicated nonzero ranges. The general and margin level bounds remain optimal after adding $j$. The coefficient $2j$ in the universal weight envelope cannot be replaced by $j$: for every $j\geq0$ there is a strictly special weight-$j$ example with $n=i=1$ whose $R^1$ boundary stalk is pure of weight $2j+2$.
\end{proposition}

\begin{proof}
For the general level bound, take the structural map of $E^n$ over $\F_5$, where $E$ is the ordinary elliptic curve of Example~\ref{c:finite-field-lie}. Its Frobenius eigenvalues have $5$-adic slopes $0,1$. In $H^i(E^n)$ the smallest slope is $\max(0,i-n)$ and the weight is $i$. Any special constituent $V_r(-t)$ has nonnegative slopes before twisting, since eigenvalues on smooth projective cohomology are algebraic integers. Hence such a lowest-slope eigenvalue requires $t\leq\max(0,i-n)$, and therefore $r=i-2t\geq\min(i,2n-i)$. The upper bound is attained. Tensoring with the constant strictly special object $H^j(E^j)$, which has a slope-zero eigenvalue, raises the required level by exactly $j$.

For the margin case let $B=E^{n-1}$ and choose a sufficiently ample projectively normal embedding. Let $Y$ be its affine cone and $X=Y\setminus\{y\}$, where $y$ is the vertex. Then $X$ is the complement of the zero section of a line bundle on the smooth $B$, so it is regular of dimension $n$ and $X\hookrightarrow Y$ is an open immersion. The blowup of the vertex gives the line-bundle resolution. Purity and proper base change identify its local cohomology with the Gysin sequence for that line bundle:
\[
 \operatorname{coker}\bigl(L:H^{i-2}(B)(-1)\to H^i(B)\bigr)
 \longrightarrow (R^if_*\Q_\ell)_{\bar y}
 \longrightarrow\ker\bigl(L:H^{i-1}(B)(-1)\to H^{i+1}(B)\bigr).
\]
The two outer terms give a short exact sequence. For $i\leq n-1$, a slope-zero class in $H^i(B)$ cannot be in the image of $L$, whose eigenvalues have slopes at least one. It gives level $i$. For $n\leq i\leq2n-1$, put $r=2n-i-1$. The primitive degree-$r$ cohomology contains a slope-zero class for the same reason. Hard Lefschetz carries it to the right-hand kernel after $i-n$ powers of $L$, and the boundary twist adds one further Tate twist. This has weight $i+1$ and slope $i-n+1$, forcing level at least $2n-i-1$. Tensoring again with $H^j(E^j)$ proves sharpness with $+j$. For $n=1$ the same construction is simply $\Gm\hookrightarrow\mathbb A^1$.

For weights with $j=0$, if $i\leq n$ use the origin-boundary stalk of
$\Gm^i\times\mathbb A^{n-i}\hookrightarrow\mathbb A^n$, whose degree-$i$ Kummer class has weight $2i$. If $n<i\leq2n$, put $q=2n-i$ and $r=i-n$. The product of the degree-$q$ class of $\Gm^q$ with the top class of $\mathbb P^r$ has degree $i$ and weight $2(q+r)=2n$. Take the product open immersion followed by projection to $\mathbb A^q$; if $q=0$ this is the projective structural map. For the margin range $i\leq2n-1$ we have $q\geq1$, so the same example satisfies the extra condition.

Finally let $U=\mathbb A^1_{\F_5}\setminus\{0,1\}$ and let $\pi:\mathcal E\to U$ be the Legendre family
\[
 y^2=x(x-1)(x-t).
\]
At $t=0$ its discriminant $16t^2(1-t)^2$ has order two and $c_4=16(1-t+t^2)$ is a unit. The node has tangent slopes whose squares are $-1$, which is a square in $\F_5$, so the reduction is split multiplicative. Thus on $H^1$ inertia acts by a nonzero rank-one unipotent operator, with invariant line $\Q_\ell$ and coinvariant line $\Q_\ell(-1)$. This is the Tate-curve calculation underlying SGA~7~I, IX, Proposition 2.2.5. Take $F_j=\operatorname{Sym}^j R^1\pi_*\Q_\ell$. The K\"unneth and Schur idempotents exhibit it as a direct summand of $R^j(\mathcal E^j/U)_*\Q_\ell$, so it has the required dimension-$j$ presentation. Its inertia coinvariants are $\Q_\ell(-j)$, and tame cohomology supplies one further twist:
\[
 (R^1u_*F_j)_{\overline0}=\Q_\ell(-j-1),
 \qquad u:U\hookrightarrow\mathbb A^1.
\]
This proves the claimed weight $2j+2$. Products with the Kummer and projective examples above attain $2\min(i,n)+2j$ for every $1\leq i\leq2n-1$. The finite matrix calculation of the symmetric-power monodromy and the displayed optimization inequalities are recorded in the accompanying program.
\end{proof}

\begin{proposition}[The top degree and its sharper weight]
\label{c:top-direct-image}
\status{Proved}
In the finite-type setting of Theorem~\ref{c:general-level-bound}, let $F$ be any lisse $\Q_\ell$-sheaf. Then $(R^if_*F)_{\bar y}=0$ for $i>2n$. If no component of local dimension $n$ maps to $y$, it vanishes already for $i>2n-1$. If $k(y)$ is finite and $F$ is pure of weight $j$, the degree-$2n$ stalk is pure of weight $2n+j$ when nonzero. This endpoint is attained by $\mathbb P^n\to\Spec\F_q$ with constant coefficient $H^j(E^j)$ for an elliptic curve $E/\F_q$. These assertions do not require the dimension-normalized presentation of $F$.
\end{proposition}

\begin{proof}
We first treat an integral source component of inherited local dimension $n$ whose image is not $y$. Compactify it as $v:X\hookrightarrow\bar X\xrightarrow{h}Y$, and pass to the strict localization at $y$. For a geometric point over $x\in\bar X_y$, let $c$ be the dimension of the strict local ring of $\bar X$ at $x$. The dimension formula gives $c+\operatorname{trdeg}_{k(y)}k(x)\leq n$. For $q>0$, Gabber's bound on the cohomological dimension of an open subset of a strictly local noetherian $c$-dimensional scheme gives
\[
 (R^qv_*F)_{\bar x}\ne0\ \Longrightarrow\ q\leq2c-1;
 \quad
 \dim\operatorname{Supp}\bigl((R^qv_*F)|_{\bar X_y}\bigr)
 \leq n-\Bigl\lceil\frac{q+1}{2}\Bigr\rceil.
\]
The exact input is \citet[Theorem 1.1, formula (1.1.1), p.~455]{c_illusie2014cd}. Its torsion statement passes to lisse $\Q_\ell$-sheaves by a stable lattice and the finite constructible stalks, whose inverse systems satisfy Mittag--Leffler. For $q=0$ the support is contained in $\bar X_y$, of dimension at most $n-1$: a fibre component of dimension $n$ would be a whole irreducible component mapping to $y$.

Proper base change and the composition spectral sequence express the desired stalk through
\[
 H^a\bigl(\bar X_{\bar y},(R^qv_*F)|_{\bar X_{\bar y}}\bigr).
\]
The cohomology of a constructible sheaf supported in dimension $s$ over an algebraically closed field vanishes above $2s$. For $q>0$ every nonzero term therefore has
\[
 a+q\leq2n-2\Bigl\lceil\frac{q+1}{2}\Bigr\rceil+q\leq2n-1;
\]
for $q=0$ it has $a\leq2n-2$. This proves the asserted vanishing on this component. Components of smaller local dimension obey the same stronger bound with their own dimension.

The irreducible components of regular $X$ are disjoint. The only possible contributions in degree $2n$ thus come from its $n$-dimensional components mapping to $y$. Such a component is a smooth variety over the finite field $k(y)$ in the weight assertion. Poincar\'e duality identifies its top cohomology with $H^0_c(X_{\bar y},F^\vee)^\vee(-n)$. On a connected nonproper component, this compactly supported section space is zero: a nonzero section of a lisse sheaf has support the entire component. On a proper component it is a subspace of a fibre of $F^\vee$, hence is pure of weight $-j$. Duality gives weight $2n+j$. Components are permuted through a finite group, which does not change weights. The case $n=0$ is the same finite-fibre argument. The projective-space example gives the stated optimality.
\end{proof}

\begin{proposition}[The lower-endpoint gate over a trait]
\label{c:lic-endpoint}
\status{Proved exact equivalence; the general mixed-characteristic vanishing is not asserted}
Let $R$ be a complete discrete valuation ring with finite residue field $k$, fraction field $K$, and $\ell\ne\operatorname{char}(k)$. Let $\mathcal Z/R$ be regular, proper, flat and generically smooth, with generic fibre $Z$. For $u:\Spec K\hookrightarrow\Spec R$, the assertion that
\[
 (u_*H^j(Z_{\bar K},\Q_\ell))_{\bar k}
       =H^j(Z_{\bar K},\Q_\ell)^{I_K}
 \quad\text{has weights at most }j
\]
is equivalent to surjectivity of the specialization map
\[
 H^j(\mathcal Z_{\bar k},\Q_\ell)
       \longrightarrow H^j(Z_{\bar K},\Q_\ell)^{I_K}.
\]
Thus, when $\dim Z=j$, improving the universal $i=0$ coefficient bound from $2j$ to $j$ includes exactly this local invariant cycle assertion.
\end{proposition}

\begin{proof}
\citet[Theorem 10.1(a), pp.~77--80 of the opened author version]{c_flachmorin2012topos} proves that the image of specialization is precisely the weight-at-most-$j$ part of the invariant space. The source has weights at most $j$, since the special fibre is proper. To check the step in that theorem, localization and Hochschild--Serre reduce surjectivity onto that weight part to the absence of weights at most $j$ in $H^{j+1}_{\mathcal Z_{\bar k}}(\mathcal Z_{R^{\rm sh}},\Q_\ell)$. Absolute purity and duality identify this supported group with $H^{2d+1-j}(\mathcal Z_{\bar k},\Q_\ell)^\vee(-d-1)$, where $d=\dim Z$; properness bounds its weights below by $j+1$. This proves the cited image assertion and both implications. The equivalence does not prove the surjectivity in general mixed characteristic; the known low-degree and equal-characteristic cases are recorded in \citet[\S3.1, p.~494]{c_esnaultwittenberg2016}.
\end{proof}

\begin{proposition}[The lower endpoint for abelian-type coefficients]
\label{c:abelian-invariant-weights}
\status{Proved}
Let $K$ be a complete discretely valued field with finite residue field of
characteristic $p$, and let $\ell\ne p$.  Suppose that $V$ is the
$\ell$-adic realization of a pure weight-$j$ motive which, after a finite
extension of $K$, is a direct summand of tensor constructions on the
$H^1$ of abelian varieties, their duals and Tate twists.  Finite-image
Artin factors are allowed.  Then the eigenvalues of geometric Frobenius
on $V^{I_K}$ have weights at most $j$.  In particular this holds for the
local realizations of pure absolute-Hodge or motivated motives of abelian
type over number fields.  The bound is attained for good reduction.
\end{proposition}

\begin{proof}
We give the weight--monodromy argument, including the tensor step.
After a finite extension we may suppose that all the finitely many
abelian generators have semistable reduction and all the Artin factors
are trivial.  Semistable reduction is SGA~7~I, IX, Theorem~3.6, printed
p.~351.  For one such abelian variety write
$U=T_\ell(A)\otimes\Q_\ell$, $V_0=U^{I_K}$, and let $W$ be the
toric part of its Tate module.  The N\'eron special fibre gives
\[
 0\longrightarrow W\longrightarrow V_0
       \longrightarrow T_\ell(B)\otimes\Q_\ell\longrightarrow0,
\]
where $B$ is the abelian quotient of the identity component.  The exact
sequence and the identification of the fixed part are IX, formula
(2.1.8), p.~328, and Proposition~2.2.5, pp.~332--333.  A torus has Tate
module of weight $-2$, and $T_\ell(B)$ has weight $-1$.  A polarization
and the orthogonality theorem IX, Theorem~2.4, p.~338, give
$W=V_0^\perp$ in the semistable case, as stated in Proposition~3.5,
p.~350.  Consequently $U/V_0$ is dual to $W$ with values in
$\Q_\ell(1)$ and has weight $0$.

Corollary~3.5.2, p.~351, says that inertia factors through the tame
$\Z_\ell(1)$ quotient and that $(g-1)^2=0$ for a generator $g$.
Thus its logarithm $N$ satisfies $N^2=0$ and $\ker N=V_0$.
The polarization is inertia invariant, so $N$ is skew adjoint and
$\operatorname{im}N=(\ker N)^\perp=W$.  It induces an isomorphism
\[
 N:U/V_0\xrightarrow{\sim}W(-1).
\]
All these statements are rational statements; no assumption on an
integral monodromy pairing is required.  The filtration
$0\subset W\subset V_0\subset U$ is therefore the monodromy
filtration centred at weight $-1$, with graded weights $-2,-1,0$.
Dualizing proves weight--monodromy, centred at $1$, for $H^1(A)$.
This derives exactly the form we use from the cited original pages
of SGA~7 \citep{c_sga7}.

For completeness, pure Weil--Deligne representations are stable under
the indicated tensor operations.  The monodromy filtration is
characterized by $NM_a\subset M_{a-2}(-1)$ and the isomorphisms
$N^a:\operatorname{gr}_a^M\xrightarrow{\sim}
\operatorname{gr}_{-a}^M(-a)$ for $a\geq0$.  On a tensor product
the convolution filtration has these properties: decompose the
nilpotent modules into $\mathfrak{sl}_2$ strings and apply the
Clebsch--Gordan decomposition to their tensor products.  Its graded
terms have the sums of the two graded weights.  Thus the centre
weights add.  Duality reverses the indices and weights, and a Tate
twist shifts the centre weight by an even integer.  An idempotent
commuting with the Weil action and $N$ preserves the unique
monodromy filtration; its image on every graded term is pure of the
same weight.  This proves stability under direct summands as well.
We apply these facts to the weight-$j$ part of the tensor presentation
of the given pure motive.  Its Weil--Deligne representation is pure
of weight $j$.

For a monodromy string centred at $j$, the vector killed by $N$ is
its lowest-weight vector.  Equivalently, $N$ is injective on every
positive-index graded term of its monodromy filtration.  Hence
$\ker N$ has weights at most $j$.  The original inertia invariants
are a Frobenius-stable subspace of $\ker N$ after the finite
extension.  If that extension has residue degree $f$, any lift of
its geometric Frobenius differs from the $f$th power of the original
Frobenius by an element of $I_K$.  These two operators therefore
agree on $V^{I_K}$.  The weight inequality is unchanged
on taking these powers, and therefore descends to $K$.
For a pure motive with good reduction the invariant space is the
whole pure weight-$j$ representation, proving the stated sharpness.
\end{proof}


The preceding results prove the requested upper-weight inequality, the stronger displayed uniform geometric-level bounds, and the manuscript's margin question in the explicit finite-type setting. They also settle the top coefficient endpoint and, for abelian-type coefficients over a trait, the lower endpoint. Proposition~\ref{c:lic-endpoint} identifies the remaining lower-endpoint problem rather than deriving it from a theorem about ordinary direct-image weights. Proposition~\ref{c:broader-coefficients} also removes the coefficient-dimension ambiguity for the general polygon and the manuscript's stated margin bound. Only the stronger reflected margin polygon retains the dimension-normalized hypothesis in our proof.

\subsection{C.5: good reduction in codimension one}

The quantifier on primes is the one in 12:140--142: set $S_\ell=\Spec R[1/\ell]$ and test $T_\ell$ at its height-one points. It does not require the ordinary $p$-adic Tate module to be unramified at a point of residue characteristic $p$.

\begin{theorem}[The positive purity range]
\label{c:vz-positive}
\status{Known, with its application proved}
Let $(R,\mathfrak m)$ be a two-dimensional regular local ring of mixed characteristic $(0,p)$ and let $U=\Spec R\setminus\{\mathfrak m\}$. Every abelian scheme over $U$ extends uniquely over $R$ if $p\notin\mathfrak m^p$. In particular the sufficient absolute ramification bound is $e\leq p-1$, including equality. More generally the same conclusion holds under the complete faithfully flat quotient condition of Vasiu--Zink Theorem 3.

If an abelian variety over $K=\operatorname{Frac}(R)$ has good reduction at every height-one prime, these conclusions apply to it. Equivalently, at each such prime one may use the prime-to-residue-characteristic unramifiedness test of the N\'eron--Ogg--Shafarevich criterion.
\end{theorem}

\begin{proof}
The first assertion is \citet[Corollary 4]{c_vasiuzink2010}. The more general condition is the existence of a faithfully flat complete regular local $R$-algebra $\widehat R$, of the same dimension and with algebraically closed residue field $k$, admitting a surjection to
\[
 W(k)[[T_1,T_2]]/(p-h),\qquad
 h\notin(p,T_1^p,T_2^p,T_1^{p-1}T_2^{p-1}).
\]
Their Theorem 3 gives extension and uniqueness under this condition. Since $e$ is the $\mathfrak m$-adic order of $p$, $p\notin\mathfrak m^p$ means $e\leq p-1$.

To apply this to a generic abelian variety, first spread it out over a nonempty open of $\Spec R$. There are only finitely many height-one points outside this open: each lies over a height-one component of its closed complement. A model over $R_{\mathfrak p}$ spreads out to a neighbourhood of $\mathfrak p$, since it and its group structure have finite presentation. The models glue by uniqueness of homomorphism extension over normal bases; a precise form is \citet[Proposition 3.3.1.5]{c_lan2021}. They give an abelian scheme on $U$. The purity theorem then applies. The equivalence of good reduction and the indicated unramifiedness is SGA~7~I, IX, Corollary 2.2.9 \citep{c_sga7}; the criterion is pointwise over the height-one discrete valuation rings, with $\ell$ invertible there. This proves the transfer required in the question, rather than silently assuming a model already existed on $U$.
\end{proof}

\begin{theorem}[An explicit abelian-surface counterexample at $p=2$]
\label{c:explicit-bad-reduction}
\status{Proved refutation of the unrestricted criterion}
Let $k=\overline\F_2$ and
\[
 R=W(k)[[u,v]]/(2-uv),\qquad K=\operatorname{Frac}(R),\qquad
 \mathfrak m=(u,v).
\]
Consider the elliptic schemes
\begin{align*}
 E_1:&\quad Y^2Z+(u+v^2)XYZ+YZ^2=X^3,\\
 E_2:&\quad Y^2Z+(v+u^2)XYZ+YZ^2=X^3,
\end{align*}
with origins $[0:1:0]$, and put $A=E_1\times_R E_2$. The sections
\[
 P_v=[-v:1:v^3]\in E_1(R),\qquad
 P_u=[-u:1:u^3]\in E_2(R)
\]
have order two on the generic fibre. Let $C_K\subset A_K$ be the order-two subgroup generated by $(P_v,P_u)$, and define the abelian surface
\[
 B_K=A_K/C_K.
\]
Then $B_K$ has good reduction at every height-one prime of $R$, hence its $T_\ell$ is unramified at every height-one point of $S_\ell=\Spec R[1/\ell]$ for every prime $\ell$, but it does not extend to an abelian scheme over $R$. The schematic closure of $C_K$ in $A$ is explicitly nonflat.
\end{theorem}

\begin{proof}
The ambient ring $W(k)[[u,v]]$ is regular of dimension three, and $2-uv$ has nonzero linear term $2$ in its maximal ideal. The quotient is regular local of dimension two, with maximal ideal $(u,v)$ and residue field $k$. It is consequently a domain of mixed characteristic; $2=uv$ has order two in $\mathfrak m$. For $E_a:Y^2Z+aXYZ+YZ^2=X^3$, the discriminant is $a^3-27$. For both displayed values of $a$ it is $1$ modulo $\mathfrak m$, hence a unit. Thus both equations define smooth projective elliptic schemes over $R$.

For $s=v,r=u$, or $s=u,r=v$, we have $rs=2$ and $a=r+s^2$. Substitution shows that $[-s:1:s^3]$ lies on $E_a$. The inverse formula sends its $Y$-coordinate to $-1+as-s^3=-1+rs=1$, so the section equals its inverse. Its generic fibre is not the origin, hence has exact order two.

Define two finite flat group schemes
\[
 G_u=\Spec R[x]/(x^2+vx),\qquad
 G_v=\Spec R[z]/(z^2+uz),
\]
with laws $x*x'=x+x'+uxx'$ and $z*z'=z+z'+vzz'$. These are integral Hopf laws: for example $1+ux$ is group-like and $(1+ux)^2=1$, by $uv=2$ and the relation. The inverse is $x\mapsto x$; direct substitution verifies the identity. The maps
\[
 G_u\longrightarrow E_1,\quad x\longmapsto[x:1:-v^2x],\qquad
 G_v\longrightarrow E_2,\quad z\longmapsto[z:1:-u^2z]
\]
are closed immersions. On the chart $Y=1$, the coordinate $X/Y$ pulls back to $x$, respectively $z$, proving the closed-immersion assertion there; their finite source makes the images closed in the entire elliptic scheme. Their images satisfy the elliptic equations: in the first case the substitution is
\[
 (-x+v-uv^2)(x^2+vx)+v^2x(uv-2)=0,
\]
and the other case interchanges $u,v$. Generically they are the subgroups $\{0,P_v\}$ and $\{0,P_u\}$. Equality of the group-law morphisms then holds over $R$ because the source and its square are flat over the domain $R$ and the target is separated. Put $H=G_u\times G_v\subset A$.

The constant order-two group may be written
\[
 D=\Spec R[t]/(t^2+t),\qquad t*t'=t+t'+uvtt'.
\]
There is a group homomorphism $\alpha:D\to H$ defined on coordinate rings by $x\mapsto vt$, $z\mapsto ut$. The polynomial relations and the coproducts agree directly. On $D(u)$ or $D(v)$ one of these two coordinate maps recovers $t$, so $\alpha_U$ is a closed immersion. These two opens cover $U$. On the generic fibre it is exactly $C_K$.

Its schematic image has coordinate algebra
\[
 R[ut,vt]\subset R[t]/(t^2+t),\qquad
 R[ut,vt]=R\oplus\mathfrak m t
\]
as an $R$-module. To verify this equality, every nonconstant monomial in $ut,vt$ reduces using $t^2=-t$ to an element of $\mathfrak m t$, while $ut,vt$ generate that summand. This algebra is torsion-free and has generic rank two, hence it is the schematic closure of $C_K$. Its closed fibre has dimension
\[
 \dim_k\bigl((R\oplus\mathfrak m t)/\mathfrak m(R\oplus\mathfrak m t)\bigr)
 =1+\dim_k(\mathfrak m/\mathfrak m^2)=3.
\]
A finite flat module over the local ring $R$ would be free of its generic rank, and would have closed-fibre dimension two. Thus the closure is not flat.

The quotient $B_U=A_U/\alpha_U(D_U)$ is an abelian scheme, since its kernel is finite locally free. Every height-one point is in $U$, which proves good reduction there. Suppose that $B_U$ extended to an abelian scheme $B/R$. The quotient isogeny $A_U\to B_U$ and its inverse up to multiplication by $2$ extend uniquely to homomorphisms $f:A\to B$ and $g:B\to A$ by the normal-base extension theorem cited above. Their composites equal multiplication by $2$, since they do so on the dense generic fibre. Thus $f$ has finite kernel, contained in $A[2]$, and is fibrewise a surjective isogeny. It is finite and flat: properness and finite fibres give finiteness, and fibrewise flatness of isogenies between abelian varieties, together with flatness over $R$, gives flatness by the fibre criterion. Therefore $\ker f$ is finite flat. Its generic fibre is $C_K$, so torsion-freeness identifies it with the schematic closure just computed. This contradicts the closed-fibre dimension three. No such $B$ exists.

All substitutions, discriminants, inverse formulas and Hopf-map polynomial identities have been checked over $\Z$ by \texttt{check\_explicit\_bad\_reduction.py}; no calculation in that certificate inverts $2$. The example is an explicit instance of the obstruction in \citet[Lemma 27 and Theorem 28(ii)]{c_vasiuzink2010}, including the prime $2$.
\end{proof}

\begin{remark}[What fails to extend]
\label{c:subgroup-versus-object}
\status{Proved clarification}
The abstract finite group $D_U$ already extends as $D$. It is its closed immersion into $A_U$, and consequently the exact sequence defining the isogeny, that cannot extend as a finite flat subgroup. Indeed finite locally free group schemes on a punctured regular local surface extend as objects by extending their coordinate bundles to reflexive modules, which are free in dimension two; multiplication and comultiplication extend by Hartogs. Exactness and closed-immersion conditions need not survive. This is the precise content of the finite-subgroup reduction in R6, rather than a failure to extend the abstract group. The source proof of Vasiu--Zink Lemma 27 reverses the Cartier-dual arrows in its printed embedding sentence; our displayed maps avoid that typing error. Lead C5 is therefore valid as an identification of the extension obstruction, while the unrestricted good-reduction claim is false.
\end{remark}

\subsection{C.6: what independence of \texorpdfstring{$\ell$}{ell} actually supplies}

\begin{proposition}[The degree-one and semistable abelian-type cases]
\label{c:ell-independence-cases}
\status{Proved in the stated cases}
Let $K$ be a number field, $v$ a finite place of residue characteristic $p$, and $\ell\ne p$.
For $M=h^1(A)$, unramifiedness of $M_\ell$ at $v$ is independent of $\ell$ and is equivalent to good reduction of $A$. It follows for every tensor construction and motivated summand of abelian varieties already having good reduction at $v$.

More generally suppose that $A/K$ has semistable reduction at $v$, that its Galois representations take values in its Mumford--Tate group $G$ over $K$, and that a motivated abelian-type object $M$ in the category it generates corresponds to a rational representation $r:G\to\operatorname{GL}(V)$. Then unramifiedness of $M_\ell$ at $v$ is independent of $\ell\ne p$. The same conclusion holds for potentially CM motives without a semistability assumption.
\end{proposition}

\begin{proof}
For the first assertion we use SGA~7~I, IX, Proposition 2.2.5 and Corollary 2.2.9. In its notation, the dimension of inertia invariants is $\mu+2\alpha$, whereas the full dimension is $2(\mu+\lambda+\alpha)$. Equality holds if and only if $\mu=\lambda=0$, which is exactly good reduction and is a geometric condition independent of $\ell$. Smooth proper base change then makes all tensor constructions unramified, and restriction to an invariant direct summand preserves the identity action.

For the semistable assertion, \citet[Theorem 1.2]{c_kisinzhou2025} supplies a common $G$-valued Weil--Deligne representation, up to $G(\C)$-conjugacy, for every prime, including $p$. Apply the rational representation $r$. In the semistable case the finite inertial type is trivial and the actual inertia action is $\exp(t_\ell(\sigma)\,dr(N))$. It is trivial exactly when $dr(N)=0$, since the tame character has infinite image and the exponential of a nonzero nilpotent cannot be identically one in characteristic zero. This vanishing is invariant under conjugation and extension of the coefficient field, proving the assertion. No descent from a ramified field extension has been used.

For potentially CM motives, the tannakian identification with the Taniyama group gives their realizations from one representation of the Weil group; the strict-compatibility statement is \citet[Remark 3, p.~262]{c_delignetaniyama1982}. On inertia the algebraic Hecke-character factors have a common finite-order part and the Artin factors are defined over a fixed number field of coefficients. Triviality of those matrices is unchanged by any embedding of that field into $\overline\Q_\ell$. Hence unramifiedness is independent of $\ell\ne p$ in this case as well.
\end{proof}

\begin{lemma}[A generator with potentially good reduction]
\label{c:potentially-good-generator}
\status{Proved}
Let $M$ be a rational motive of abelian type over a number field $K$, in either
the absolute-Hodge or motivated category.  Let $v$ have residue characteristic
$p$.  If $M_{\ell_0}$ is potentially unramified at $v$ for one $\ell_0\ne p$,
then, after a finite Galois extension $E/K$, the object $M_E$ belongs to the
tensor category generated by finitely many abelian varieties over $E$ which
have good reduction at every place above $v$.
\end{lemma}

\begin{proof}
We first work over $\overline\Q$, using the chosen complex embedding.
The weak Mumford--Tate lifting construction of
\citet[Theorem 2.12 and Corollary 2.15, pp.~15--20]{c_noot2006} gives an
abelian variety $B$ with a central isogeny
\[
 \pi^{\mathrm{der}}:\operatorname{MT}(B)^{\mathrm{der}}
       \longrightarrow\operatorname{MT}(M)^{\mathrm{der}}
\]
identifying their adjoint Hodge homomorphisms.  Here descent to
$\overline\Q$ is part of the argument, not an assumption: start with an
abelian generator $A/\overline\Q$ of $M$, take its Mumford--Tate decomposed
weak lift $B_0$, and descend $B_0$ by Noot's Proposition 3.1 and Corollary 3.2,
pp.~20--21.  The proof of Corollary 2.15 selects the factors of $B_0$ whose
$\Q$-simple derived groups survive in $\operatorname{MT}(M)$.  These factors
descend up to isogeny because their rational endomorphism idempotents descend
under an algebraically closed field extension.  Rational motives only require
this isogeny-level descent.

The proof of Noot's Proposition 3.4, pp.~22--23, applies with the Hodge
realization of $M$ in place of $H^1(A)$ and gives
\begin{equation}
\label{c:weak-lift-generates}
 M_{\overline\Q}\ \in\
 \langle h^1(B),\text{CM motives}\rangle^{\otimes}.
\end{equation}
To make this extension of the source statement explicit, use the common
group formed from the central tori and $\operatorname{MT}(B)^{\mathrm{der}}$.
Its representation on $\omega_B(M)$ is generated by the faithful
representation on $H^1(B)$ and representations of its torus quotient.
For an irreducible constituent $W_2$ of $\omega_B(M)$, choose an irreducible
constituent $W_1$ of a tensor construction on $H^1(B)$ and a torus
representation $W_3$ with $W_2\subset W_3\otimes W_1$.
The nonzero contraction $W_2\otimes W_1^\vee\to W_3$ shows that $W_3$ itself
is an actual polarizable Hodge subquotient, of CM type, with integral
Hodge degrees.  Full faithfulness of Betti realization for abelian motives
and the realization of CM Hodge structures by CM abelian varieties prove
\eqref{c:weak-lift-generates}.  In the motivated category the same step uses
Andr\'e's Theorem 0.6.2.  Only finitely many CM varieties and morphisms are
needed for this particular $M$, so all these data descend to a finite
extension of $K$.

Put $H=\operatorname{MT}(M_{\overline\Q}\oplus h^1(B))$.
Its two projections have the same adjoint Hodge datum.  The Hodge
homomorphism defining $H$ lies in the fibre product over that adjoint group,
and projects onto both Mumford--Tate groups.  The kernel of the projection
$H\to\operatorname{MT}(M)$ lies in $\{1\}\times Z(\operatorname{MT}(B))$,
so is central in $H$ and meets its semisimple derived group in a finite
group.  Thus $H^{\mathrm{der}}\to\operatorname{MT}(M)^{\mathrm{der}}$ is a
central isogeny.  After a finite extension, the arithmetic realizations of
$M\oplus h^1(B)$ take values in $H(\Q_\ell)$ for every $\ell$: finitely many
absolute Hodge tensors cut out $H$, and one extension defines all of them.

Fix a place above $v$ and take a further finite local extension over which
$M_{\ell_0}$ is unramified and $B$ is semistable.  An inertia element acts on
$M_{\ell_0}\oplus H^1(B,\Q_{\ell_0})$ as
\[
 u=(1,\rho_{B,\ell_0}(\tau))\in H(\Q_{\ell_0}),
\]
and $u$ is unipotent by the semistable monodromy theorem.  Its image in the
torus $H/H^{\mathrm{der}}$ is trivial, so $u\in H^{\mathrm{der}}$.
Its image in $\operatorname{MT}(M)^{\mathrm{der}}$ is trivial, and hence
$u$ lies in a finite kernel.  A finite-order unipotent in characteristic
zero is the identity.  Thus $H^1(B,\Q_{\ell_0})$ is unramified, and
N\'eron--Ogg--Shafarevich gives good reduction of $B$ over this local
extension.  The same reasoning applies at \emph{every} place above $v$,
since potential unramifiedness of $M_{\ell_0}$ transfers to all those
completions.

The auxiliary CM abelian varieties are potentially good as well: after
fixing their Hodge tensors and taking a semistable extension, their inertia
is both unipotent and contained in a torus, hence trivial.  The same
criterion applies.  There are only finitely many varieties and places.
Choose finite extensions of the completions giving good reduction and
realize them within one finite global extension by weak approximation and
Krasner's lemma.  Pass to its normal closure over $K$; good reduction is
preserved by further extension.  This gives the asserted $E$ and proves
the direction of transfer required in the lemma.
\end{proof}

\begin{theorem}[Independence of unramifiedness for all abelian-type motives]
\label{c:ell-independence-abelian}
\status{Proved, using the October 2026 rationality theorem}
Let $M$ be a rational absolute-Hodge or motivated motive of abelian type
over a number field $K$, and let $v$ have residue characteristic $p$.
If $M_{\ell_0}$ is unramified at $v$ for one $\ell_0\ne p$, then
$M_\ell$ is unramified at $v$ for every $\ell\ne p$.
The statement includes $p=2$, and includes $\ell=2$ when $p\ne2$.
\end{theorem}

\begin{proof}
Apply Lemma~\ref{c:potentially-good-generator}.  Over a finite Galois
extension $E/K$, the motive $M_E$ is a summand of a finite direct sum $P$
of motives $h^{i_a}(A_a)(m_a)$, where the $A_a$ are powers or products of
abelian varieties with good reduction at \emph{every} place above $v$.
Indeed, tensor constructions on $h^1$, their duals and Tate twists occur
as K\"unneth summands in such powers.  A finite direct sum permits different
degrees $i_a$ and twists $m_a$.

Restriction and induction give morphisms
\[
 M\longrightarrow\operatorname{Ind}_E^K(M_E)\longrightarrow M
\]
whose composite is $[E:K]$.  After division by this nonzero rational
integer, $M$ is a summand of $\operatorname{Ind}_E^K(P)$.
These are finite \'etale correspondence maps in either motivic category;
after extension to $E$ they are the diagonal and summation maps over
$\Gal(E/K)$.  Thus the descent assertion is established over the original
$K$, without inferring it from trivial inertia over a ramified extension.

Write $e$ for the resulting rational Hodge, equivalently motivated,
projector.  The induced ambient varieties are finite disjoint unions of
the conjugate $A_a$.  Fix the chosen place of $\overline\Q$ above $v$ and
take a common finite Galois local extension $L/K_v$ splitting the finite
extensions of induction and giving good models for all components.
Smooth proper base change shows, simultaneously for every $\ell\ne p$,
that inertia on the ambient object factors through the same finite
quotient
\[
 \Delta=I_v/(I_v\cap G_L).
\]
It follows that this also holds on $M_\ell$.  Tate twists are trivial on
this inertia group when $\ell\ne p$.

By uniqueness of the abelian scheme models, the descent maps extend to
them.  On the geometric special fibres, each $\tau\in\Delta$ acts by a
permutation of components and origin-preserving algebraic isomorphisms
$u:X_0\to Y_0$ between conjugate abelian varieties.  Inertia fixes the
algebraic residue field, so these maps are actual homomorphisms over that
field.  Their pullbacks describe the cohomological inertia action.

The graph of every such $u$ is a rational Lefschetz class on $X_0\times Y_0$.
If $d=\dim Y_0=\dim X_0$ and $f(x,y)=y-u(x)$, then in every Weil
cohomology,
\begin{equation}
\label{c:inertia-graph-lefschetz}
 [\Gamma_u]=f^*[0_{Y_0}]
   =f^*\left(\frac{\theta_{Y_0}^{d}}
                         {\deg(\theta_{Y_0}^{d})}\right),
\end{equation}
where $\theta_{Y_0}$ is any ample divisor class.  This uses the common
degree normalization of top classes.  The expression on the right is one
rational polynomial in divisor classes, with the same cycle realizations
for all $\ell$.

Use the correspondence convention
$Z_*(x)=(\mathrm{pr}_2)_*([Z]\smile\mathrm{pr}_1^*x)$.
Then $\Gamma_u$ acts by $u_*=(u^{-1})^*$, and the alternating trace of
$e_*\circ u^*$ is the intersection pairing of $e$ with $\Gamma_u$.
For a projector supported in degree $i$, the ordinary trace is therefore
$(-1)^i$ times that pairing.  On a disjoint union, the same identity pairs
the matrix entries of $e$ with the graphs indexed by the component
permutation.  For a direct sum of different degrees and twists, only
the diagonal blocks relative to those Galois-stable summands contribute
to the trace.  Consequently
\[
 \operatorname{Tr}(\tau\mid M_\ell)
\]
is a finite sum, with the indicated signs, of pairings of specialized
rational Hodge classes with complementary rational Lefschetz classes on
products of abelian varieties with good reduction.  The generic matrix
entries of $e$ are Hodge classes on these products; realization and smooth
proper specialization preserve their actions, tensor operations, duality
and K\"unneth decompositions.

The rationality theorem \citep[Theorem 1.1]{c_rationality2026}, also
\citep[Theorem 1.1, pp.~1--2]{c_rationalityexposition2026}, says that every
such pairing is one rational number in every $\ell$-adic realization
$\ell\ne p$, including residue characteristic $2$.  Hence for each
$\tau\in\Delta$ there is a common rational trace $t_\tau$.
The assumed unramifiedness at $\ell_0$ gives
$t_\tau=\dim_\Q\omega_B(M)$.  In every other realization the operator
$\tau$ has finite order.  Its eigenvalues are roots of unity, and under a
complex embedding their sum equals their number.  Equality in the triangle
inequality forces every eigenvalue to be $1$.  Finite-order operators in
characteristic zero are semisimple, so $\tau$ is the identity on $M_\ell$.
This holds for every $\tau$, proving the theorem.
\end{proof}

\begin{remark}[The exact role of the new input]
\label{c:ell-independence-new-input}
\status{Proved scope clarification}
The proof does not apply N\'eron--Ogg--Shafarevich to an arbitrary original
generator: such a generator need not have good reduction when its summand
is unramified.  Lemma~\ref{c:potentially-good-generator} changes the
generator and proves its potential good reduction.  The finite descent
action is then recovered over $K$ by the common trace pairing.  Since
\eqref{c:inertia-graph-lefschetz} is Lefschetz, the proof uses the
rationality theorem alone; it does not use the CM Hodge theorem or assume
algebraicity of the Hodge projector.  It proves independence of
unramifiedness and, for potentially unramified objects, independence of
the finite inertia character.  It does not assert full Weil--Deligne
compatibility for objects with nonzero monodromy.
\end{remark}

% Supplementary reduction for C.6.  This does not settle the general iff.
\begin{proposition}[A purity reduction for homogeneous Chow motives]
\label{c:chow-unramified-purity-reduction}
\status{Proved conditional implication}
Let $K$ be a number field, $v$ a finite place with residue field of
cardinality $q$, and $p$ its residue characteristic.  Let
$M=(X,e,m)$ be a rational Chow motive over $K$, where $X$ is smooth
projective of dimension $d$ and the rational algebraic projector
$e\in\CH^d(X\times_KX)_{\Q}$ acts only in cohomological degree $i$.
Suppose that, for one prime $\ell_0\ne p$, the representation
$M_{\ell_0}$ is unramified at $v$ and the eigenvalues of geometric
Frobenius on it are algebraic numbers all of whose complex conjugates
have one common absolute value.  Then $M_\ell$ is unramified at $v$
for every $\ell\ne p$.

In particular, the asserted independence follows for such a motive if
the weight--monodromy conjecture holds for
$H^i(X_{\overline{K_v}},\Q_{\ell_0})$: its unramified direct summand
$M_{\ell_0}$ is then pure of weight $i-2m$.
\end{proposition}
\begin{proof}
Write $V_\ell=eH^i(X_{\overline{K_v}},\Q_\ell)(m)$.
Saito's Theorem~0.1, in its more general form Theorem~3.1,
gives a rational alternating trace independent of $\ell$ for the
composition of an algebraic correspondence and any Weil element
\citep[Theorems~0.1 and 3.1]{c_saito2003weights}.
The theorem extends from integral to rational correspondences by
linearity.  Taking the correspondence $e$, its action on all other
degrees is zero; since it is idempotent and commutes with the Galois
action, one obtains the exact equality
\[
 \operatorname{Tr}(\sigma\mid V_\ell)
 =(-1)^i q^{-m n(\sigma)}
   \sum_a(-1)^a\operatorname{Tr}
   \bigl(e\circ\sigma\mid H^a(X_{\overline{K_v}},\Q_\ell)\bigr).
\]
Thus $\operatorname{Tr}(\sigma\mid V_\ell)$ is a common rational
number for each $\sigma\in W_{K_v}$.  Apply this also to all positive
powers of $\sigma$.  Newton identities show that the characteristic
polynomial of every such $\sigma$ is in $\Q[T]$ and independent of
$\ell$.  In particular the dimension $r$ is common; if $r=0$ there
is nothing to prove.

Fix a lift $\phi\in W_{K_v}$ of geometric Frobenius.  The local
monodromy theorem supplies, for every $\ell\ne p$, a nilpotent operator
$N_\ell$ and an open inertia subgroup on which the action is
$\exp(t_\ell N_\ell)$; it also gives
\[
 F_\ell N_\ell F_\ell^{-1}=q^{-1}N_\ell,
 \qquad F_\ell=\rho_\ell(\phi).
\]
The common characteristic polynomial of $F_\ell$ has roots of the
single absolute value in the hypothesis.  Over an algebraic closure
of the coefficient field, $N_\ell$ maps a generalized
$\lambda$-eigenspace of $F_\ell$ into its generalized
$q^{-1}\lambda$-eigenspace.  The latter eigenvalue cannot occur:
under any embedding of the common algebraic splitting field into
$\C$, multiplication by $q^{-1}$ changes its absolute value.
Consequently $N_\ell=0$ for every $\ell$.

The image of inertia in each $\rho_\ell$ is therefore finite.  For
each $\tau\in I_{K_v}$ the common trace is
$\operatorname{Tr}(\rho_{\ell_0}(\tau))=r$.  A finite-order matrix
in characteristic zero is semisimple with roots of unity as
eigenvalues.  Their sum can equal the positive integer $r$ only if
all of them are $1$, as follows by embedding their cyclotomic field
into $\C$ and using equality in the triangle inequality.  Hence
$\rho_\ell(\tau)=1$ for every $\tau$ and every $\ell\ne p$.

For the last assertion, the weight--monodromy filtration is functorial
for the algebraic projector.  A pure Weil--Deligne direct summand
whose monodromy is zero is concentrated in monodromy degree zero,
so its Frobenius roots all have weight $i-2m$.
\end{proof}

\begin{remark}[Exact scope of this reduction]
\label{c:chow-purity-scope}
\status{Proved scope clarification; the general claim remains unresolved}
For the full group $H^i(X)$, the purity assumption in the proposition
is Serre's unramified-purity conjecture C7
\citep[\S2.3, printed p.~19-07]{c_serre1970local}.
For a prescribed proper summand it is the corresponding
\emph{motivic extension} of that condition; the literal original C7
must not be silently enlarged.  The named weight--monodromy
conjecture for the ambient cohomology is a sufficient alternative.
The use of Saito's theorem requires an algebraic correspondence.
It cannot be applied to an arbitrary absolute-Hodge or motivated
projector without an additional comparison result.  No converse to
any of these named conjectures is proved here, and this proposition
does not close the general-motive requirement in C.6.
\end{remark}

\begin{remark}[The remaining general-motive claim]
\label{c:ell-independence-gate}
\status{Unresolved in the general motivic category}
Theorem~\ref{c:ell-independence-abelian} settles the whole abelian-type assertion, including its finite-inertia descent. The direct argument from an arbitrary original generator would still be invalid: $\bigwedge^2H^1(E)=\Q_\ell(-1)$ is unramified at a place of residue characteristic different from $\ell$, even when $E$ has bad reduction. The weak-lift lemma and the trace argument are what repair that inference.

For general motives, the precise named sufficient conjecture is Fontaine's motivic Weil--Deligne compatibility conjecture: \citet[\S2.4.3, Conjecture $C_{\mathrm{WD}}(X,m)$, p.~342, and its motivic formulation in \S2.4.7, pp.~344--345]{c_fontaine1994pst}. The motivic formulation is required to retain the specified projector; compatibility of an unlabelled ambient cohomology representation alone does not identify every prescribed summand. Write the local Weil--Deligne data as $(r_\ell,N_\ell)$, with $r_\ell$ of finite inertial image. On an open subgroup of inertia the actual action is $\exp(t_\ell(\tau)N_\ell)$, where the tame character has infinite image. Consequently
\[
 M_\ell\text{ is unramified at }v
 \quad\Longleftrightarrow\quad
 N_\ell=0\ \text{ and }\ r_\ell|_{I_v}=1.
\]
Both conditions are invariant under extension of coefficient fields and isomorphism of Weil--Deligne representations. Thus Fontaine's compatibility, which only requires isomorphism after common coefficient extension and not a chosen rational form, implies the requested independence. His weaker monodromy-graded character formulation in \S2.4.5 also suffices: the characters at the identity determine whether the monodromy filtration is concentrated in degree zero, equivalently $N=0$; after that, the finite inertia character detects triviality.

This proves an implication, not the exact equivalence requested in the question. Algebraically, the representations trivial on a subgroup $J\subset G$ detect its closed normal closure, not the whole inertia representation. Even all tensor tests do not formally recover the nilpotent's Jordan type: in $\mathfrak{sl}_3$, both $E_{12}$ and $E_{12}+E_{23}$ act nontrivially in every nontrivial algebraic representation, since the differential's kernel is an ideal in the simple Lie algebra. Their ranks differ. The accompanying notes give compatible abstract Weil--Deligne pairs and a matrix certificate; no such pair is asserted to come from a global motive, and it is not a motivic counterexample.

The theorem of \citet[Theorem 5.1]{c_commelin2019} concerns abelian motives and quasi-compatibility; by itself it does not supply finite-inertia traces for an arbitrary motive at a bad place. Similarly the $G$-valued expectation of \citet[Conjecture 1.1]{c_kisinzhou2025} is stronger than the asserted triviality test. A proof in the general motivic category, or an exact equivalence with a published named conjecture of precisely the required strength, remains an outstanding assertion, recorded separately from the now proved abelian-type theorem.
\end{remark}
